\section{Partial Truth and Fixed Points}
\label{sec:p11-ktf-proofs}

\begin{annexitem}{annex:p11.ktf.effective}{thm:p11-ktf-effective}{Hyperarithmetical total definability}
\begin{proof}
For either specified scheme, the total-extension upper theorem shows
that the positive extension of a formula defined at all numerical
inputs is lightface \(\Delta^1_1\). The independently constructed
Theorem~\ref{thm:p04-doh-characterization} identifies exactly this
pointclass with the sets computable in a native jump oracle \(H_a\)
for \(a\in\mathcal O\). Thus every such total extension belongs
to \(\mathsf{HYP}\).

Conversely suppose \(A\leq_T H_a\) for a valid native notation
\(a\). The other inclusion of that same arithmetic characterization
gives \(A\in\Delta^1_1\). The alternating product-tree
construction in the total-converse theorem supplies an actual
monadic formula, defined at every input, whose positive extension
is \(A\) under Strong Kleene and under the original two-coordinate
unrestricted supervaluation. This proves the reverse inclusion.
The arithmetic graph proof imported by the jump construction is
independent of the desired characterization; neither this argument
nor its premise identifies HYP with a new definition of \(\Delta^1_1\).
There is no arbitrary set parameter or change of the ordinal
notation system in either inclusion.
\end{proof}
\end{annexitem}

\begin{annexitem}{annex:p11.ktf.pi-operations}{thm:p11-ktf-pi-operations}{Effective positive operations on analytical definitions}
\begin{proof}
Write a \(\Pi^1_1\) predicate as \(\forall X\,\eta(\vec n,X,Z)\)
with an arithmetical matrix and optional fixed free set parameter \(Z\).
Finite conjunction pairs its universal set variables. Finite disjunction
uses
\[
 (\forall X\,\eta(X))\lor(\forall Y\,\xi(Y))
 \quad\Longleftrightarrow\quad
 \forall X\forall Y\,(\eta(X)\lor\xi(Y)).
\]
If neither disjunct holds, its two counterexamples refute the right
side; the other direction is immediate. Universal number quantification
commutes with universal set quantification. For existential number
quantification the required equivalence is
\[
 \exists n\forall X\,\eta(n,X,Z)
 \quad\Longleftrightarrow\quad
 \forall Y\exists n\,\eta(n,Y_n,Z),
\]
where \(Y_n\) is the \(n\)-th set coded by numerical pairing in
\(Y\). A witness on the left works against every array. If the
left fails, each \(n\) has an arithmetical counterexample set;
the arithmetical-choice lemma supplies one array of all counterexamples,
refuting the right side. This does not interchange an arbitrary
existential and universal quantifier without that argument.
Each displayed transformation is effective on finite formula codes.
Computable substitution uses the unique arithmetic graph of its
substitution function, as in the preimage argument for the upper bound.

Induct on a positive arithmetical formula built from set-membership
atoms, truth-free arithmetical atoms, Boolean conjunction/disjunction
and number quantifiers. Replacing membership atoms by their
\(\Pi^1_1\) definitions preserves the class by the preceding
transformations. Negative truth-free atoms remain arithmetical;
no negation of the replaced predicate is admitted. If universal set
quantifiers are additionally prefixed, combine them with the generated
universal prefix, retaining their occurrences as parameters in the
arithmetical matrix transformations. Fresh variable renaming prevents
capture. This establishes the stated effective and uniform closures.
\end{proof}
\end{annexitem}

\begin{annexitem}{annex:p11.ktf.pi-comparison}{thm:p11-ktf-pi-comparison}{A recursive-ordinal analytical comparison}
\begin{proof}
First justify the universal index. Recognizing the second-order
definition syntax is computable. Translate its first-order numerical
matrix, after numerical instantiation, into the language with one
predicate interpreted by its set argument. Its code is computed from
\((e,z)\). The relation \(\operatorname{Val}(R,V)\) in the upper
bound proof verifies its classical truth set arithmetically. Thus the
assertion that every set argument and its unique valuation make the
matrix true has a universal set prefix and arithmetical matrix.
This constructs \(\mathsf U\) uniformly, without a first-order
uniform arithmetic truth predicate. The effective tree normal form
gives its recursive family \(\mathcal S_{e,z}\).

The Kleene--Brouwer relation is a computable linear order on a recursive
tree: two distinct finite sequences either have a proper-prefix
relation or a first unequal entry. Comparison at that first entry,
with a prefix placed above its extensions, is transitive as well.
An infinite branch gives an infinite descending sequence of its
successive prefixes. Conversely suppose there is an infinite
descending Kleene--Brouwer sequence. The empty node occurs at most
once. After it, the first entries cannot increase and hence eventually
stabilize. A node consisting of that one entry can then occur at most
once, and the second entries eventually stabilize for the same reason.
Inductively every position stabilizes, and every sufficiently late
node has the corresponding stabilized prefix. Prefix closure puts
all of these prefixes in the tree, yielding a branch. At each position
the eventual value and its least stabilization stage are uniquely
specified numerical choices. Thus no branch means that the recursive
linear order is a well-ordering, and its order type \(\rho\) is
a recursive ordinal.

A coded function is an initial-segment isomorphism from the second
order to the first when it maps their numerical fields appropriately,
is injective, preserves and reflects their order, and has downward
closed image. These conditions quantify only over numbers and the
function graph. Extend its values outside the source field by \(0\)
if a total numerical function code is used. For two well-orders such
an isomorphism exists exactly when the source order type is at most
the target order type. This follows from their unique ordinal
representations: take the isomorphism to the corresponding initial
segment, or transport an existing one to the ordinals. An ill-founded
linear order cannot embed in a well-order, since the image of its
descending sequence would contradict well-foundedness.

Now \(\operatorname{Below}(p,q)\) is the conjunction of a
\(\Pi^1_1\) well-foundedness assertion and the universal-set
assertion that no proposed graph is the stated isomorphism. It is
therefore \(\Pi^1_1\). If \(q\) is well founded it compares
the two ordinal types strictly, with \(p\) required to be well
founded. If \(q\) is ill founded, every well-founded \(p\)
passes the no-isomorphism test and no ill-founded \(p\) passes
the first conjunct. These are exactly the claimed cases.
\end{proof}
\end{annexitem}

\begin{annexitem}{annex:p11.ktf.effective-induction}{thm:p11-ktf-effective-induction}{An effective monotone-induction bound}
\begin{proof}
For numerical parameters \(m,z\), define
\[
 S_{m,z}=\{u:\operatorname{Below}((m,u),(m,z))\},\qquad
 Q(m,z)\quad\Longleftrightarrow\quad z\in\Gamma(S_{m,z}).
\]
The comparison is \(\Pi^1_1\) uniformly. The assumed compiler
applied to its set definition, followed by identifying the output
argument with \(z\), supplies a \(\Pi^1_1\) definition of
\(Q(m,z)\). We next construct an actual self-index for this
definition; an equation naming the desired set alone would not suffice.

Let \(\operatorname{diag}_{\Pi}(k)\) be the total computable
code operation which, on a \(\Pi^1_1\) formula with two designated
free number variables, substitutes \(\bar k\) for the first and
leaves the second free. On invalid inputs return the code of a fixed
false \(\Pi^1_1\) formula. Its arithmetic graph is available.
Form the \(\Pi^1_1\) formula
\(H(k,z)=Q(\operatorname{diag}_{\Pi}(k),z)\), replacing the
informal graph application and putting its prefix in the specified
syntax by the effective-operation lemma. In its arithmetical matrix
retain the tautological conjuncts \(k=k\) and \(z=z\); these keep
both designated free variables even if the operator's output is a
constant predicate. Let \(h\) be this actual
formula code, and put \(m^*=\operatorname{diag}_{\Pi}(h)\).
The universal-index property and the substitution computation give
\[
 \mathsf U(m^*,z)\quad\Longleftrightarrow\quad
 H(h,z)\quad\Longleftrightarrow\quad Q(m^*,z).
\]
Define \(A=\{z:\mathsf U(m^*,z)\}\). This has not assumed
that \(A\) is the desired inductively generated set.

If \(z\notin A\), the comparison lemma gives
\(S_{m^*,z}=A\). The self-index equality then says
\(z\notin\Gamma(A)\). Hence \(\Gamma(A)\subseteq A\):
\(A\) is pre-fixed. Starting with the empty set, the increasing
bottom iteration of \(\Gamma\) is below \(A\) at every stage,
by monotonicity and union limits.

For the reverse inclusion use transfinite induction on the ordinal
\(\rho(m^*,z)\) for \(z\in A\). Its comparison set is precisely
the members of \(A\) with smaller \(\rho\). By induction each
such member \(u\) is in the bottom iteration by stage
\(\rho(m^*,u)+1\), which is at most \(\rho(m^*,z)\).
Thus \(S_{m^*,z}\) is included in that stage. The self-index
equation says \(z\in\Gamma(S_{m^*,z})\), so monotonicity puts
\(z\) into stage \(\rho(m^*,z)+1\). The initial case uses
the empty smaller-rank set, and a limit rank uses the same union
inclusion; no uniform collection of stages is presumed.

Each \(\rho(m^*,z)\) is the order type of a recursive well-order.
Its successor is recursive too, by adjoining one largest element to
a computable presentation. Therefore each stage just obtained is
strictly below \(\omega_1^{\mathrm{CK}}\). The union stage at
that ordinal contains \(A\), and was already shown below \(A\),
so equals it. This union stage is below its own jump, as is every
stage of a monotone bottom iteration; pre-fixedness gives the reverse
inclusion. It is consequently a fixed point. Stage induction makes
it below every fixed point, and all later stages remain equal.
This proves both assertions and supplies the effective upper bound
without an unproved admissibility premise.
\end{proof}
\end{annexitem}

\begin{annexitem}{annex:p11.ktf.ck-exact}{thm:p11-ktf-ck-exact}{Exact arithmetic closure ordinal}
\begin{proof}
For Strong Kleene use a numerical signed-rule operator
\(\Omega\): retain the valid signed-sentence portion of the input,
add the atomic and non-code
truth-atom base signs, and add the conclusions of every applicable
Boolean, quantified, equal-denotation and truth-prefix closure rule
defining \(W\). Malformed or non-sign inputs are discarded, so
every output lies in the effective set of signed sentences. Its
fixed sets are thus exactly the typed closed sets defining \(W\).
Its defining formula is arithmetical with only
positive occurrences of membership in the input set. In particular
the all-instances rules use universal number quantifiers; they do
not test nonmembership of the input. It is monotone and its effective compiler
on \(\Pi^1_1\) definitions is supplied by the operations lemma.
Effective induction therefore gives its least fixed
set \(W\) by the CK stage, with every signed sentence entering
at a strictly earlier stage.

Compare its bottom stages with the original truth iteration.
Induct that a sign in \(\Omega^\alpha\) is a defined value at
\(K^\alpha\). For atomic base signs this is immediate; the
negative non-code truth atoms have their value from stage \(1\).
Boolean and quantified forcing conclusions have their value whenever
their premise signs have their values, hence at the preceding truth
stage and by persistence at the next one. Equal-denotation replacement
preserves atomic values. A truth-prefix rule advances at most one
truth stage, because evaluating \(A\) at stage \(\alpha\)
puts its code into the appropriate coordinate at \(\alpha+1\).
Retained signs persist. At limits, each sign comes from an earlier
stage. This proves the comparison for all ordinals.

Consequently every sign in \(W\) is evaluated at some truth stage
\(\alpha<\omega_1^{\mathrm{CK}}\). Its sentence code is present
by \(\alpha+1<\omega_1^{\mathrm{CK}}\). The CK truth union
therefore has all the positive and negative codes of the least fixed
point; it also has all non-codes negative since stage \(1\). It is
below that point by leastness and thus equals it. This proves the
Strong-Kleene upper bound.

For supervaluation use the ambient signed operator \(G\) from the
upper-definability proof. For a signed sentence its defining
condition is, explicitly, universal quantification over \(R,V\)
of the disjunction of
\[
 \exists k\,(k\in X_1\land k\notin R),\quad
 \exists k\,(k\in X_2\land k\in R),\quad
 \neg\operatorname{Val}(R,V),\quad
 \text{the indicated classical sign in }V.
\]
There are only positive occurrences of the input coordinates.
The negative non-code output is arithmetical and independent of
the input. The effective-operations lemma therefore converts every
\(\Pi^1_1\) input definition into a \(\Pi^1_1\) output
definition, uniformly with numerical parameters. The ambient
operator is monotone, so effective induction bounds its bottom
iteration by CK. That bottom iteration is exactly the consistent
two-coordinate supervaluation iteration already established; there
is no substitution of a positive-only scheme or a different seed.

Finally the cofinal recursive-order tree construction gives the
CK lower bound for each of these same two iterations. Combining
the upper and lower bounds proves the two exact equalities.
\end{proof}
\end{annexitem}

\begin{annexitem}{annex:p11.ktf.total-converse}{thm:p11-ktf-total-converse}{Total definition of every lightface delta-one-one set}
\begin{proof}
For \(A\in\Delta^1_1\), apply the normal-form lemma separately
to its \(\Pi^1_1\) definition and to that of its complement.
Let their recursive trees be \(\mathcal T_n^0,\mathcal T_n^1\),
so exactly one of the two root trees is well founded at each \(n\).
Apply the diagonal product construction to these actual uniform
trees. For a pair of valid nodes, the immediate pair descendants
advance both coordinates by one edge. Whenever one coordinate
subtree is well founded, this pair relation is well founded: an
infinite pair path would project to a branch in that coordinate.
Equivalently a failure of well-foundedness on numerical pair codes
would yield such a path by successively taking the least available
pair descendant. This ensures that well-founded induction applies.

At either least fixed point, use that induction to prove every
pair sentence in this well-founded pair relation defined. Assume
all its immediate pair sentences defined. Their code bits are
then bivalent by fixed-point transparency. The ordinary arithmetic
guards and unique code graphs make the displayed sentence's
equation exactly
\[
 v(\vartheta_{s,t})=
   \text{the value of }\forall a\in\operatorname{Child}(s)
                   \exists b\in\operatorname{Child}(t)
                        v(\vartheta_{a,b}).
\]
All active instances here are bivalent. Inactive guards are
classically false and make their conjunctions false or their
implications true even if irrelevant truth atoms have gaps.
Thus the whole formula has a defined Strong-Kleene value.
For supervaluation every extending total interpretation agrees
on all the active defined code bits, so the same classical value
is unanimous. This proves totality by well-founded induction,
including the empty-child cases: an empty outer universal is
true, while an empty inner existential is false.

If the first subtree is well founded and the second has a branch,
the pair value is true. Prove this by induction on the first
subtree, while keeping a branch in the second. For each first
child use the next child on that branch; the induction hypothesis
gives a true pair value, which witnesses its inner existential.
Every first child therefore satisfies it, proving the universal.
If the first subtree has a branch and the second is well founded,
the pair value is false. Induct on the second subtree, keep the
next first-tree child on its branch, and apply the induction
hypothesis to every second child. All its inner instances are
false, so its existential is false and witnesses failure of the
outer universal. The two induction arguments include their
respective leaf cases and hold under both schemes by the
bivalent equation just established.

Let \(g(n)\) be the computable code of the sentence for the
pair of roots at \(n\), and let its graph be \(G_A(x,y)\).
At either least point the formula
\(\exists y(G_A(x,y)\land T(y))\) reads exactly that root's
defined bit. It is true on \(A\) and false on its complement,
and defined at every input. The reverse inclusion is the already
proved total-extension upper bound. This yields the exact
\(\Delta^1_1\) class rather than a Boolean combination of
the two gap-valued weak representations.
\end{proof}
\end{annexitem}

\begin{annexitem}{annex:p11.ktf.evaluation}{thm:p11-ktf-evaluation}{Evaluation and persistence}
\begin{proof}
Terms receive unique denotations by induction on their finite formation:
a variable denotes its assigned object, a constant its fixed interpretation,
and an application of a function uses the already determined denotations
of its arguments. An induction on terms also shows that their denotations
depend only on the variables occurring in them. Each formula has a finite
formation tree. Its atomic case therefore has a unique value, and the
Boolean and quantifier clauses assign a unique value from the values
already assigned to proper subformulas and their assignments. For a
quantifier the true and false conditions cannot both hold: a universal
cannot have all instances true and one false, and an existential cannot
have one instance true and all false. Thus the residual case is unambiguously
undefined. This proves existence and uniqueness by formula induction.

For coincidence, suppose assignments agree on the free variables of
\(A\). The atomic case follows from term coincidence, and Boolean
cases follow from their operands. For \(\forall x B\) or \(\exists x B\),
changing both assignments to send \(x\) to the same \(d\) makes them
agree on all free variables of \(B\). Apply the induction hypothesis
to every \(d\in D\) and then the quantifier clause. A sentence has
no free variables, so its value is assignment independent.

The term-substitution identity is
\(u[t/x]^{\mathcal M,g}=u^{\mathcal M,g[x:=t^{\mathcal M,g}]}\),
proved by term induction. Formula induction gives the corresponding
identity for atomic formulas and Boolean constructors. For a quantifier
binding \(x\), no free occurrence under its binder is replaced, and
the quantifier evaluation overrides the changed assignment. For a binder
\(y\ne x\), choose the bound variable fresh from the variables of
\(t\), as required for capture-free substitution. Then
\(t^{\mathcal M,g[y:=d]}=t^{\mathcal M,g}\), and changing \(x\)
and \(y\) commutes. Apply the induction hypothesis to the operand
at every \(d\), and then its quantifier clause. Renaming a bound variable
to a fresh one preserves evaluation by the same assignment argument.
This establishes the substitution claim in its capture-free form.

For persistence, ordinary atomic values are unchanged. If a truth atom
is true at \(S\), its denotation belongs to \(S_1\subseteq R_1\),
and remains true at \(R\); the false case uses the second coordinate.
An undefined atom can take any later value without violating
\(\leq_i\). Negation preserves the same relation by exchanging
the two defined values. For conjunction, a true value has two true operands,
which remain true; a false value has a false operand, which remains false.
For disjunction, a true operand remains true, and two false operands remain
false. Undefined compound values impose no additional restriction.
For a true universal, every instance remains true; for a false universal,
its false instance persists. For a true existential, its true witness
persists; for a false existential, all instances remain false. Formula
induction proves the stated information-order inequality.
\end{proof}
\end{annexitem}

\begin{annexitem}{annex:p11.ktf.arithmetic-graphs}{thm:p11-ktf-arithmetic-graphs}{Arithmetic formulas for computable numerical graphs}
\begin{proof}
The full finite-sequence, pairing, configuration and successful-history
construction is Lemma~\ref{thm:p04-doh-arithmetic-graphs}, with proof
\ref{annex:p04.doh.arithmetic-graphs}. It gives the arithmetic graph of
each partial computable numerical function effectively from its machine.
For a computable predicate use its total characteristic function and
require graph output one. The chosen finite-word and substitution
operations are computable numerical functions under the specified codes,
so the same construction applies to them. A total function has its unique
actual output at each standard-model input. No truth predicate or
hyperarithmetical characterization is assumed in that earlier derivation.
\end{proof}
\end{annexitem}

\begin{annexitem}{annex:p11.ktf.tree-normal-form}{thm:p11-ktf-tree-normal-form}{Well-founded-tree normal form and arithmetic choice}
\begin{proof}
Consider first \(\exists X\,\eta(n,X,Z)\), with arithmetical
matrix and optional fixed oracle \(Z\). Replace the set by its
characteristic function and require its values to be \(0\) or \(1\).
Put the number quantifiers of the matrix in prenex form by ordinary
classical renaming and Boolean transformations. For each existential
number quantifier introduce a function of its preceding universal number
arguments. These functions can be chosen when the prefix is true:
at each argument tuple choose the least natural witness leaving the
remaining suffix true. Conversely any functions satisfying the resulting
universal matrix supply witnesses for the original prefix, by induction
on its finite sequence of quantifiers. This numerical least-witness
selection uses no choice principle for arbitrary sets.

Pack the characteristic function and the finitely many witness functions
into one function \(g:\mathbb N\to\mathbb N\), with numerical
pairing coding a component and its argument tuple. The assertion is
equivalent to existence of \(g\) satisfying a uniformly specified
countable list of tests. These tests require binary characteristic
values and the quantifier-free arithmetic matrix at each tuple of
universal numbers. Each test is an oracle computation using finitely
many values of \(g\) and of \(Z\), and halts when all its queries
are available. Form a tree of finite sequences \(s\) by the following
algorithm: reject \(s\) if one of the first \(|s|\) tests, run for
\(|s|\) steps using the entries of \(s\), finishes and reports
failure. A query to an unavailable entry is left unfinished. Include
the empty sequence. This tree is recursive relative to \(Z\).
Its membership is prefix closed: every rejection of a shorter sequence
is also detected for an extension with the greater time and test bound.

A successful full \(g\) has all prefixes in the tree. Conversely, if
\(g\) is a branch, a failing test would halt after finitely many
steps using finitely many entries. A sufficiently long prefix would
be rejected. Hence all tests hold, giving the required witness. This
proves the existential-set normal form. Negating it gives the asserted
recursive well-founded-tree representation for a universal-set formula
with no oracle parameter. The algorithm for the tree is obtained
effectively from the given finite formula.

For the uniform choice assertion, apply this construction to each
\(\exists Y\,\theta(j,Y,Z)\). Each resulting tree has a branch.
A nonempty set of branches of a tree on natural numbers has a canonically
chosen branch: start at the empty sequence and repeatedly choose the
least next natural number for which the extended prefix still has a
branch. Such a number exists at every step; the resulting infinite
sequence is a branch because the tree is prefix closed. The property
of having a branch is used as a mathematical predicate here, not as
an effective decision procedure. Apply this unique prescription at
each \(j\). Separation and Replacement collect the chosen branches,
and extract their characteristic-set components. Pairing the index
\(j\) with its entries yields one set coding all the witnesses.
No effective bound or hyperarithmetical selection is asserted.
\end{proof}
\end{annexitem}

\begin{annexitem}{annex:p11.ktf.pi-upper}{thm:p11-ktf-pi-upper}{Upper definability bound for arithmetic truth}
\begin{proof}
Encode a signed sentence by pairing its effective code with its sign.
The assertion \(\operatorname{Closed}_{\mathrm{SK}}(X)\) that a set
of signed codes satisfies the closure conditions defining \(W\) is
arithmetical relative to \(X\). Atomic values and equal-denotation
tests are computable, and the finitely many constructor and substitution
operations have arithmetic graphs. Boolean forcing rules quantify only
over their numerical codes. A universal-premise rule has the form
\(\forall k\,\langle c(A[\bar k/x]),i\rangle\in X\)
implying membership of its numerical conclusion. Thus even its infinite
premise family adds only number quantifiers, not set quantifiers.
By the signed-closure theorem,
\[
 \langle c(A),i\rangle\in W
 \quad\Longleftrightarrow\quad
 \forall X\,[\operatorname{Closed}_{\mathrm{SK}}(X)
                  \to\langle c(A),i\rangle\in X].
\]
This is a lightface \(\Pi^1_1\) definition of both signed extensions.

For supervaluation use the full powerset of signed domain elements,
including inconsistent subsets. If \(X\) is such a set, let
\(X_1=\{n:\langle n,1\rangle\in X\}\) and
\(X_2=\{n:\langle n,0\rangle\in X\}\). Define a numerical
operator \(G(X)\) by the original supervaluation jump: a sentence
code receives sign \(i\) when every total interpretation extending
these coordinates makes it have that sign; non-codes receive only
the negative sign. This gives a monotone operation on a full powerset,
even if an inconsistent input makes the completion family empty.
Its iteration from the empty set is exactly the previously proved
consistent supervaluation iteration, so its least fixed point there
is the signed encoding \(U\) of \(K_{\mathrm{sv}}\).
No lattice property of consistent partial states is being assumed.

Define \(\operatorname{Val}(R,V)\) to mean that \(V\) codes the
classical true closed sentences when \(T\) is interpreted by the
set \(R\). This relation is arithmetical in \(R,V\): demand no
invalid sentence codes, the correct computed values for atomic
equalities and truth atoms, both directions of each Boolean clause,
and \(\forall x A\in V\) iff all numerical closed instances are
in \(V\), with the analogous existential clause. Its atomic
term-denotation and instance-code operations have arithmetic graphs.
For each \(R\) there exists exactly one such \(V\). Existence
is the ordinary classical recursive satisfaction set; uniqueness
follows by induction on the number of logical constructors, unchanged
by substituting a closed numeral into an operand. This uses a set of
truth values as an extra set variable, and does not claim a uniform
first-order definition of arithmetic truth.

The condition \(G(X)\subseteq X\) now says two things. All negative
non-code signs belong to \(X\). For every missing signed sentence
code \(j\), there exist \(R,V\) with
\(X_1\subseteq R\), \(R\cap X_2=\varnothing\),
\(\operatorname{Val}(R,V)\), and the opposite classical value
for the sentence at \(j\). These witness conditions are arithmetical
in \(j,X,R,V\). For an index already in \(X\), impose no
condition on its witnesses. The arithmetical-choice lemma therefore
packs all countermodel and valuation witnesses into one set \(Y\).
Consequently there is an arithmetical relation \(\operatorname{Pre}(X,Y)\)
such that
\[
 G(X)\subseteq X\quad\Longleftrightarrow\quad
 \exists Y\,\operatorname{Pre}(X,Y).
\]

An increasing bottom stage is below every pre-fixed set \(X\) with
\(G(X)\subseteq X\): prove this from the empty stage by monotonicity
at successors and unions at limits. Since \(U\) is itself a fixed
point, membership in \(U\) is equivalent to membership in all
pre-fixed sets. Hence
\[
 j\in U\quad\Longleftrightarrow\quad
 \forall X\forall Y\,[\operatorname{Pre}(X,Y)\to j\in X],
\]
a lightface \(\Pi^1_1\) definition after pairing its universal
set variables. Empty completion families at arbitrary inconsistent
\(X\) cause no semantic contradiction in this ambient numerical
operator; the actual bottom iteration and its fixed point are consistent,
as already proved.

Finally substitution of a numeral into a fixed monadic formula is a
total computable map \(f(n)=c(A[\bar n/x])\). The preimage of a
\(\Pi^1_1\) set under such a map is \(\Pi^1_1\): replace
\(\forall X\,\eta(f(n),X)\) by
\(\forall X\exists y\,[\operatorname{Graph}_f(n,y)\land\eta(y,X)]\).
The output is unique and independent of \(X\), and the remaining
quantifiers are numerical. Transparency at either fixed point identifies
these code preimages with the true and false extensions of the formula.
\end{proof}
\end{annexitem}

\begin{annexitem}{annex:p11.ktf.tree-ranks}{thm:p11-ktf-tree-ranks}{Tree height as a truth rank}
\begin{proof}
The numerical word operation defining \(F\) is total and computable,
so the graph lemma provides its truth-free arithmetic formula. Forming
\(B\), coding it and substituting its numeral gives the displayed
sentences. The retained parameter equality makes different node
numerals occur in different strings, proving injectivity. At any
partial state, a Strong-Kleene value of \(\psi_z\) is therefore
the universal conjunction of the membership values of exactly its
child codes: false guards make all other numerical instances true,
and the unique graph output for a genuine child reads its code.
It is true iff all child codes are positive, false iff some child
code is negative, and undefined otherwise.

The same criterion holds for unrestricted supervaluation. If every
child code is positive, every completion makes the universal true.
If one is not positive, choose a completion assigning it false,
which witnesses failure of supertruth. If a child is negative the
universal is false in every completion. If no child is negative,
the total state with positive coordinate equal to the given positive
coordinate together with all child codes is a compatible completion
making the universal true. Thus superfalsity fails in that case.
These arguments also cover no children, when the sentence is true.

Starting from the empty coordinates, induct that no family code is
negative at any stage: the preceding criterion could create a
negative sign only from an already negative child code; a limit
cannot introduce the first one. Absence of a branch gives a genuinely
well-founded immediate-child relation here: if a nonempty subset of coded
nodes had no minimal member, repeatedly choosing its least-numbered
available child would give a branch. This numerical prescription needs
no choice principle. For a well-founded node, use well-founded
induction on its children. If their first evaluation
ranks are \(\alpha_y\), their codes first enter at
\(\alpha_y+1\). The node's sentence first becomes true at
\(\sup_y(\alpha_y+1)\): at that stage all those codes are present,
and at any lower stage one required code is absent. An empty
child family gives rank \(0\). The recursion is exactly the height
recursion, proving the rank equality for each scheme.

If there is a branch through \(z\) and its sentence were true at
some stage, let its first evaluation rank be \(\alpha\). The
next node on the branch would have its code positive by that stage,
so would have a strictly smaller first evaluation rank. Continuing
along the given branch produces a strictly decreasing infinite
sequence of ordinals, impossible: the set of its entries would
have a least one followed by a still smaller entry. No such truth
value can occur. Since false values were excluded, the sentence
and its code remain undefined throughout.
\end{proof}
\end{annexitem}

\begin{annexitem}{annex:p11.ktf.pi-representation}{thm:p11-ktf-pi-representation}{Positive monadic extensions and analytical completeness}
\begin{proof}
For a given lightface \(\Pi^1_1\) set \(A\), obtain its recursive
tree family from the normal-form lemma and its actual sentences
\(\psi_z\) from the diagonal forest construction. Let \(r(n)\)
be the computable node code of the empty root of \(\mathcal T_n\),
and put \(g(n)=c(\psi_{r(n)})\). The tree-rank theorem gives
\[
 n\in A\quad\Longleftrightarrow\quad g(n)
 \text{ belongs to the positive least-point extension},
\]
under either scheme. This is the required computable many-one
reduction. The upper bound already proved makes each positive
code extension \(\Pi^1_1\)-complete.

Let \(G_A(x,y)\) be the truth-free arithmetic graph of \(g\).
The monadic formula
\(\exists y(G_A(x,y)\land T(y))\) has positive extension exactly
\(A\). Under Strong Kleene, all graph-false instances are false,
and the unique graph-true instance has the value of the selected
truth atom. Under supervaluation, its classical value in every
completion is likewise the bit assigned to that unique code.
At the fixed point that bit is true precisely on \(A\) by the
root construction. Every other positive monadic extension is
\(\Pi^1_1\) by the upper theorem, proving equality of the classes.
\end{proof}
\end{annexitem}

\begin{annexitem}{annex:p11.ktf.ck-lower}{thm:p11-ktf-ck-lower}{Cofinal arithmetic ranks and the Church--Kleene lower bound}
\begin{proof}
Choose a computable well-ordering \(\prec\) of a computable
subset \(E\subseteq\mathbb N\) with order type \(\alpha\).
Its tree consists of finite strictly decreasing sequences from
\(E\), together with the empty root. This is recursive and has
no branch. A node whose last entry is \(e\) has height equal
to the order type of \(\{d\in E:d\prec e\}\), by induction
along the well-ordering. The root height is therefore
the supremum of these predecessor order types plus one, which
is \(\alpha\); for an empty order it is \(0\).
The tree sentence for that root has precisely this evaluation
rank under both schemes. Its code is not entered until
\(\alpha+1\), so an iteration stable earlier cannot be the
bottom iteration: all stages after a stable stage are equal.
Each closure ordinal is at least \(\alpha+1\) for every
recursive \(\alpha\). Taking their supremum gives the stated
\(\omega_1^{\mathrm{CK}}\) lower bound. This establishes
cofinality, not the separate effective upper bound.
\end{proof}
\end{annexitem}

\begin{annexitem}{annex:p11.ktf.total-upper}{thm:p11-ktf-total-upper}{Total monadic extensions are lightface delta-one-one}
\begin{proof}
Let \(A\) be the formula's true extension and \(B\) its false
extension at the least point. Each is \(\Pi^1_1\) by the upper
theorem. Total definedness and the absence of inconsistent values
give \(B=\mathbb N\setminus A\). Therefore \(A\) and its
complement are both \(\Pi^1_1\), equivalently \(A\) is both
\(\Pi^1_1\) and \(\Sigma^1_1\). This is exactly the asserted
lightface \(\Delta^1_1\) conclusion. It does not construct a
total formula for each such set or prove its jump-hierarchy
characterization.
\end{proof}
\end{annexitem}

\begin{annexitem}{annex:p11.ktf.finite-equivalence}{thm:p11-ktf-finite-equivalence}{Effective equivalence of the two finite hierarchies}
\begin{proof}
All code operations in this proof manipulate finite words with a
computably specified finite auxiliary vocabulary. They have truth-free
arithmetic graphs by the graph lemma. Extend each sentence translation
to a total numerical map by sending non-sentence inputs of its source
language to \(0\), which is not a sentence code in the target language.
Replace an informal application of such a code map inside an object
formula by existential quantification over its unique arithmetic graph.
Choose its new bound variables fresh. These conventions make every
translation below an actual formula in the stated target language.

First show how to lift lower noncumulative predicates. Let
\(a_i(n)=c(\mathsf N_i(\bar n))\), with this code taken in
\(L_N^i\). This map is total and computable. Since \(\mathsf N_{i+1}\)
is truth for that language, the formula
\[
 \exists y\,(\operatorname{Graph}_{a_i}(x,y)
                         \land\mathsf N_{i+1}(y))
\]
has exactly the value of \(\mathsf N_i(x)\) for every numerical
input. This includes non-sentence inputs, since the atom named by
\(a_i(n)\) is always a legitimate sentence and its predicate then
has its stipulated false value. Replacing each \(\mathsf N_i\)-atom
by this formula lifts formulas of \(L_N^i\) to \(L_N^{i+1}\).
Induction on formulas proves preservation for every assignment: the
atomic fact was proved, Boolean constructors are unchanged, and each
quantifier ranges over the same natural numbers. Repeating the lift
gives a computable value-preserving embedding from \(L_N^i\) into
\(L_N^r\) for every \(i<r\).

Now construct the two desired translations together by induction on
\(r\), starting with the identity at level \(0\). Suppose the
lower-level maps have been constructed. For \(u_r\), replace
\(\mathsf N_r(t)\) by
\[
 \exists y\,(\operatorname{Graph}_{u_{r-1}}(t,y)
                                      \land\mathsf C_r(y)).
\]
If the denotation of \(t\) is a sentence code of \(L_N^{r-1}\),
the induction hypothesis gives a true-value-equivalent sentence in
\(L_C^{r-1}\). Its truth value is precisely the value read by
\(\mathsf C_r\), so the displayed formula equals the old
\(\mathsf N_r\)-atom. On every other input the graph returns
\(0\), and both predicates are false. This proves atomic preservation,
and formula induction again proves preservation under all assignments.

For the reverse translation, first define the current predicate
\(\mathsf C_r\) in \(L_N^r\) by
\[
 \exists y\,(\operatorname{Graph}_{v_{r-1}}(x,y)
                                      \land\mathsf N_r(y)).
\]
The same valid-code and default-code argument proves its values correct.
Each earlier \(\mathsf C_i\), with \(i<r\), already has a defining
formula in \(L_N^i\) by the construction at that earlier level.
Lift that formula into \(L_N^r\) by the noncumulative embeddings
proved in the first part. Substitute these respective defining formulas
for all \(\mathsf C_i\)-atoms to construct \(v_r\). Their arithmetic
graph formulas use only already constructed lower-level translations;
there is no recursive request for the map currently being defined.
The syntactic replacements are finite, computable, capture free and
preserve the original free variables. The atomic equalities and another
formula induction establish the stated preservation property.

Both translations therefore preserve every formula's numerical extension.
They do not have to undo one another as strings: the graph quantifiers
and predicate-definition expansions can remain in their composites.
The argument supplies the finite arithmetic equivalence and makes no
claim about an unspecified transfinite language-level convention.
\end{proof}
\end{annexitem}

\begin{annexitem}{annex:p11.ktf.order}{thm:p11-ktf-order}{Chain unions and a monotone jump}
\begin{proof}
Reflexivity and transitivity follow from inclusion, and mutual inclusion
of both coordinates implies equality of states. Empty coordinates are
below every state. For a pairwise compatible family \(\mathcal A\),
suppose an element belonged to both coordinates of its union. It would
belong positively to some \(S\in\mathcal A\) and negatively to some
\(R\in\mathcal A\), contrary to compatibility. Thus the union is
a state. It is an upper bound, and any upper bound contains both of its
coordinates, establishing leastness. Comparable states are compatible,
because a conflict in the smaller state and the larger one would be a
conflict in the larger state itself. This proves the chain assertion.

In the jump, a sentence cannot have both values, and the injective coding
prevents distinct sentences from assigning opposing signs to one code.
The negative non-codes are disjoint from every sentence code. The output
coordinates are therefore disjoint. Persistence gives both coordinate
inclusions when \(S\preceq R\); the non-code portion is identical.
Finally take any \(d\in D\). The states \((\{d\},\varnothing)\)
and \((\varnothing,\{d\})\) have no upper bound in the space of
consistent states. Since \(D\) is nonempty, this proves that the full
space is not a lattice.
\end{proof}
\end{annexitem}

\begin{annexitem}{annex:p11.ktf.construction}{thm:p11-ktf-construction}{Least fixed extension and countable closure}
\begin{proof}
First justify the recursion without treating an ordinal-length definition
as a proof of existence. For any set ordinal \(\gamma\), a candidate
initial segment is a function on an ordinal \(\delta\leq\gamma+1\)
satisfying the initial, successor and limit clauses at each index in its
domain. Two such functions agree on their common domain: otherwise their
least disagreement is not the initial index, and its successor or limit
clause is determined by agreeing earlier values, a contradiction. Use
transfinite induction to extend these unique agreeing initial segments.
At a successor, append the uniquely prescribed jump. At a limit,
Replacement collects the already unique shorter segments, their union
is a function on the limit, and the coordinate union prescribes the next
value. All collected functions and states are sets. This constructs the
unique recursion through \(\gamma\); constructions for different
\(\gamma\) agree by the preceding uniqueness argument.

Induct simultaneously that all values are states, that the sequence is
increasing, and that each state is below its jump. The initial below-jump
assertion is the seed hypothesis. If it holds at \(\alpha\), monotonicity
gives \(S^{\alpha+1}=\Phi(S^\alpha)\preceq
\Phi(S^{\alpha+1})\). Consistency of the successor follows from the
disjoint-output lemma. At a limit \(\lambda\), the earlier values are
a chain, so their union is a state and contains each one. Each earlier
\(S^\beta\) is below \(S^{\beta+1}\preceq\Phi(S^\lambda)\).
Taking unions shows \(S^\lambda\preceq\Phi(S^\lambda)\).
These facts establish all the induction assertions, including comparison
between any two indices.

From stage \(1\) onward, every non-code is permanently in the negative
coordinate, and none is in the positive coordinate. Thus a strict
successor increase at any \(\alpha\geq1\) adds a signed sentence
code. Fix an injective index for \(C\times\{1,0\}\) into
\(\mathbb N\), obtained from the supplied countable listing by taking
the first occurrence of each code and its sign. Suppose there were no
closure stage below \(\omega_1\). For every \(\xi<\omega_1\),
the step from \(S^{\xi+1}\) to \(S^{\xi+2}\) would then add at
least one signed sentence code. Choose the added code of least numerical
index. Different steps select different codes, since the sequence is
increasing and a code, once assigned, persists. This would inject
\(\omega_1\) into \(\mathbb N\), contradicting its definition.
The selection is uniquely specified by a numerical least element and
uses no choice of enumerations or witnesses. Existence of \(\omega_1\)
in ZF follows from Hartogs's construction applied to \(\mathbb N\).
Consequently some \(\theta<\omega_1\) is a closure stage, and the
least such stage exists by the well-ordering of the ordinals.

At that stage \(S^\theta=\Phi(S^\theta)\), so it is a fixed
point. Successors preserve it, and a limit union of copies of it is again
it, proving stability at all subsequent stages. Let \(F\) be any
fixed point with \(S\preceq F\). Induction gives
\(S^\alpha\preceq F\): the initial case is assumed, successors
use monotonicity and \(\Phi(F)=F\), and limits use coordinate unions.
Thus \(S^\theta\preceq F\), proving leastness among fixed extensions
of the seed. For the bottom seed the hypothesis \(S\preceq F\)
holds for every fixed point, giving the least fixed point \(K\).
\end{proof}
\end{annexitem}

\begin{annexitem}{annex:p11.ktf.examples}{thm:p11-ktf-examples}{Grounded values and the self-referential cases}
\begin{proof}
Truth-free atoms have their classical values even at the bottom state.
The Boolean and quantifier clauses restrict to their classical clauses
when all relevant operand and instance values are defined. Structural
induction gives a defined value for every truth-free sentence at stage
\(0\), hence rank \(0\).

If \(A\) first has sign \(i\) at stage \(\alpha\), the jump puts
its code into coordinate \(i\) at \(\alpha+1\). There was no
earlier entry: a first entry cannot occur at a nonzero limit union, since
membership there comes from an earlier stage; an earlier successor entry
would have come from an earlier defined evaluation of \(A\). Stage
\(0\) contains no codes. Persistence also rules out any earlier
opposite sign. This proves the precise offset between evaluation rank
and code entry.

For the liar, write \(t\) for the value assigned to its code by \(T\)
at a fixed point. The actual code computation gives
\(v_F(\lambda)=\neg_i t\), where \(\neg_i\) exchanges \(0\)
and \(1\) and fixes \(\mathbf u\). Fixed-point transparency also
gives \(v_F(\lambda)=t\). Neither defined value equals its negation,
so \(t=\mathbf u\). The same reasoning works at every fixed point.

At the bottom state the truth-teller's code is unassigned. A successor
assigns it just when the preceding evaluation of \(\tau\) was defined,
but that evaluation is exactly its preceding membership value; a limit
cannot introduce it for the first time. Transfinite induction therefore
leaves it undefined throughout the bottom sequence.
On the other hand the state
\((\{\ulcorner\tau\urcorner\},\varnothing)\) is inflationary:
\(\tau\) evaluates true there and its code is returned positively.
Similarly \((\varnothing,\{\ulcorner\tau\urcorner\})\) is
inflationary because \(\tau\) evaluates false. Their least fixed
extensions supplied by the construction theorem respectively make
\(\tau\) true and false. This also proves that the two points have
no compatible common extension.

Finally a name of the code of \(A\) evaluates under \(T\) to the
sign of that code, which at a fixed point is precisely the evaluated
sign of \(A\). If their common value is \(0\) or \(1\), both
directions of the material biconditional are true. If it is undefined,
the negations and disjunctions forming that biconditional remain
undefined by the displayed tables.
\end{proof}
\end{annexitem}

\begin{annexitem}{annex:p11.ktf.limit-example}{thm:p11-ktf-limit-example}{A sentence settled at the first limit}
\begin{proof}
Every \(A_0\) has value true at every stage. The bottom state assigns
no sentence code, so \(A_k\) with \(k>0\) is undefined there.
For a positive finite stage \(j+1\), the truth atom \(A_{k+1}\)
reads the code of \(A_k\); its value is the value of \(A_k\) at
stage \(j\). Induction on \(j\), simultaneously for all \(k\),
therefore gives truth exactly for \(k\leq j\), and no false value.
At finite stage \(j\), the instance of \(\sigma\) with input
\(k=j\) reads the code of \(A_j\), which first enters at stage
\(j+1\), so that instance remains undefined. No instance is false.
Consequently \(\sigma\) is undefined at every finite stage.

Each code \(q(k)\) has entered positively by stage \(k+1\).
Their union at \(\omega\) therefore makes every instance
\(T(q(\bar k))\) true. Since the domain is standard \(\mathbb N\),
these are all the instances of its universal quantifier. The value of
\(\sigma\) at \(K^\omega\) is true, and its first code entry
is at \(\omega+1\). The coordinate union of the finite jumps equals
\(K^\omega\), which does not yet contain that code; its jump does.
This proves the strict inequality and shows why monotonicity alone does
not imply preservation of increasing countable unions.
\end{proof}
\end{annexitem}

\begin{annexitem}{annex:p11.ktf.maximal}{thm:p11-ktf-maximal}{Maximal extensions}
\begin{proof}
Fix a fixed point \(F\), and partially order its fixed extensions by
\(\preceq\). This is a nonempty set, containing \(F\). If a chain
of such extensions is nonempty, let \(U\) be its coordinate union.
It is a state extending \(F\). Each member \(G\) satisfies
\(G=\Phi(G)\preceq\Phi(U)\), so \(U\preceq\Phi(U)\).
Its closure \(\operatorname{Cl}(U)\) is consequently a fixed upper
bound of that chain. For the empty chain use \(F\). Zorn's principle
now gives a maximal fixed extension of \(F\). It is maximal among
all fixed points as well, since any fixed point extending it also
extends \(F\).

Suppose a fixed point \(G\) leaves \(\tau\) undefined. Add its
code to the positive coordinate to obtain \(U\). This is consistent.
Monotonicity gives \(G=\Phi(G)\preceq\Phi(U)\), and the newly
added code is also positive in \(\Phi(U)\), since \(\tau\) reads
that very code. Hence \(U\preceq\Phi(U)\). Its closure is a proper
fixed extension of \(G\). One can instead add the negative sign and
obtain another proper extension by the same argument. A maximal fixed
point thus decides the truth-teller. Every fixed point leaves the liar
undefined, as already proved, so maximality does not imply totality.
The two fixed points deciding \(\tau\) oppositely have no common
consistent upper bound; a greatest fixed point is therefore impossible.
The compound \(\kappa\) is undefined at the least point, since both
\(\tau\) and \(\lambda\) are undefined there. At any fixed point
deciding \(\tau\) true, its first disjunct makes \(\kappa\) true.
At a point deciding \(\tau\) false, its first disjunct is false and
its second is the conjunction of true with the undefined liar, so
\(\kappa\) is undefined. The two points deciding \(\tau\)
oppositely can each be extended to a maximal point by the argument above.
Their decided truth-teller values persist. They give maximal points
with the two claimed behaviours of \(\kappa\).
\end{proof}
\end{annexitem}

\begin{annexitem}{annex:p11.ktf.intrinsic}{thm:p11-ktf-intrinsic}{The lattice of intrinsic fixed points}
\begin{proof}
We first establish the compatibility property used for closure. If
\(S\) is compatible with a fixed point \(F\), their coordinate
union \(U\) is a state. Monotonicity puts \(\Phi(S)\) and
\(\Phi(F)=F\) below the consistent state \(\Phi(U)\); hence
\(\Phi(S)\) is compatible with \(F\). If an inflationary seed
\(S\) is compatible with \(F\), transfinite induction now makes
every stage of its iteration compatible with \(F\): use this argument
at successors, and the absence of a conflicting element in any earlier
coordinate at a limit union. Its closure is therefore compatible with
\(F\) as well.

The least fixed point \(K\) is below every fixed point and is intrinsic.
Let \(\mathcal A\) be any nonempty family of intrinsic fixed points.
They are pairwise compatible, so their union \(U\) is a state.
Each \(G\in\mathcal A\) is below \(\Phi(U)\), since
\(G=\Phi(G)\). Taking unions gives \(U\preceq\Phi(U)\).
The union is compatible with every fixed point, because each of its
members is. The closure compatibility argument therefore shows that
\(\operatorname{Cl}(U)\) is intrinsic. It is a fixed upper bound
of \(\mathcal A\), and is below every fixed upper bound by leastness
of the closure. It is thus the join of that family within intrinsic
fixed points. For the empty family the union is the bottom state and
its closure is \(K\), giving the empty join as well.

Intrinsic fixed points form a set, being a definable subset of the state
space. Take their join to obtain a greatest intrinsic fixed point.
For any family \(\mathcal B\) of intrinsic fixed points, its intrinsic
lower bounds form a set and contain \(K\) if the family is nonempty.
Their join is below every member of \(\mathcal B\), since each member
is a fixed upper bound of all those lower bounds. Every intrinsic lower
bound is below that join. It is therefore the greatest lower bound.
For empty \(\mathcal B\), every intrinsic fixed point is a lower
bound and their join is the greatest intrinsic point, as required for
the empty meet. This proves complete-lattice existence on this restricted
set without applying a lattice theorem to the full partial-state space.
\end{proof}
\end{annexitem}

\begin{annexitem}{annex:p11.ktf.intrinsic-example}{thm:p11-ktf-intrinsic-example}{Intrinsic truth beyond groundedness}
\begin{proof}
For a fixed point \(F\), write \(t=v_F(\rho)\). The two truth
atoms in its defining string have values \(t\) and \(\neg_i t\)
by transparency. Hence \(t=t\lor_i\neg_i t\). This equality is
impossible for \(t=0\), possible for \(t=1\), and possible for
\(t=\mathbf u\). Thus \(\rho\) is false at no fixed point.
At the bottom its own code and its negation's code are both unassigned.
Simultaneously induct that neither code acquires a value: their sentence
values remain undefined at a successor, and neither can first enter at
a limit. Consequently \(\rho\) is ungrounded.

Let \(S=(\{\ulcorner\rho\urcorner\},\varnothing)\). Here the
first disjunct of \(\rho\) is true, so \(S\preceq\Phi(S)\).
It is compatible with every fixed point because none makes \(\rho\)
false. Its closure is therefore intrinsic by the compatibility argument
in the preceding proof and still makes \(\rho\) true. This proves
its intrinsic value.

For \(\chi=\tau\lor\neg\tau\), the tables make its value
\(1\) when \(\tau\) is defined, and \(\mathbf u\) otherwise.
There is no false case. If an intrinsic fixed point gave \(\chi\)
a defined value, it would decide \(\tau\). A true decision conflicts
with the fixed point making \(\tau\) false, and a false decision
conflicts with the one making it true. Both are impossible for an
intrinsic point. Thus \(\chi\) has no intrinsic value, despite
never being false among fixed points.
\end{proof}
\end{annexitem}

\begin{annexitem}{annex:p11.ktf.boundaries}{thm:p11-ktf-boundaries}{Three diagnostic boundaries}
\begin{proof}
Every state \(S\) extends to the total state
\((S_1,D\setminus S_1)\); the second coordinate contains \(S_2\)
by disjointness. Classical evaluation makes a sentence or its negation
true in every total interpretation, so \(\lambda\lor\neg\lambda\)
is supertrue at every partial state. At a Strong-Kleene fixed point
the liar and its negation are both undefined, and their disjunction is
undefined by the table.

For a classical total revision, let \(b_j\in\{0,1\}\) be the
membership bit of the liar code at stage \(j\). Its current sentence
value is \(1-b_j\), because its actual string negates the truth
atom naming its own code. The next interpretation assigns that evaluated
value to its code, so \(b_{j+1}=1-b_j\). The two successive total
states differ on this code and cannot extend one another in the information
order. This calculation proves the successor oscillation and does not
choose a policy for revision at limit stages.

Finally the fixed-point liar proof excludes its code from both
\(F_1\) and \(F_2\). Its nonmembership in \(F_1\) is therefore
a true assertion in the classical metalanguage. The object-language
sentence \(\neg T(d(\bar b))\), which is the liar itself, is
undefined under the stipulated truth and negation clauses. These two
assertions use different predicates and are not equivalent descriptions
within the object language.
\end{proof}
\end{annexitem}

\begin{annexitem}{annex:p11.ktf.empirical}{thm:p11-ktf-empirical}{One syntax, three semantic outcomes}
\begin{proof}
When \(P\) is empty, every instance of its implication has a false
antecedent; the Strong-Kleene disjunction clause makes it true regardless
of the value of the truth atom. Its universal is true already at the
bottom state. When \(P\) contains exactly the code of \(B\),
the sole nontrivial instance is \(\neg T(\ulcorner B\urcorner)\),
with the code understood as the denotation of the quantified variable.
At the bottom it is undefined, and all other instances are true; thus
\(H\) is undefined. The truth-free sentence \(B\) is true at
stage \(0\), so its code belongs to the positive extension at
stage \(1\). The nontrivial instance is then false, making \(H\)
false of rank \(1\). Finally, when \(P\) contains exactly the
code of \(H\), evaluation of \(H\) in any state is the negation
of the value assigned to its own code, all other instances being true.
At a fixed point this value must equal its own negation, which rules out
both defined values. This is paradoxicality with unchanged sentence syntax.
For the mutual-dependence case, both \(H_1\) and \(H_2\) are
undefined at stage \(0\), because their selected truth atoms are
undefined and neither universal has a false instance. At stage \(1\)
the code of the true base sentence \(B\) is positive. Its occurrence
under negated truth makes one instance of \(H_2\) false, settling
\(H_2\) false. The code of \(H_2\) is still unassigned there,
so \(H_1\) remains undefined. At stage \(2\) that code is negative,
and the sole nontrivial instance of \(H_1\) becomes true. All its
other guarded instances are true as well, giving rank \(2\). This
calculation depends on the given ordinary predicate extensions, not on
any preassigned ordering of the two sentences.
\end{proof}
\end{annexitem}

\begin{annexitem}{annex:p11.ktf.tarski}{thm:p11-ktf-tarski}{Finite hierarchy embedding}
\begin{proof}
The syntactic recursion is on the numerical index of the guard, not on
truth. The fixed base vocabulary determines its sentence strings and
\(C_1\). By the syntactic-adequacy hypothesis it supplies a truth-free
formula \(A_1\) defining that set. Once \(A_r\) is fixed, substitution
of its finite formula for the auxiliary token determines \(C_{r+1}\)
syntactically, and the same hypothesis supplies its next defining guard.
At no step is a truth value used to decide membership in a guard set.
Every guard, being truth-free, has its classical bivalent value from the
initial stage and at every later stage. The proof below uses precisely
this stated adequacy; it does not infer syntactic definability from
the mere existence of an external injection of codes.

Every sentence with code in \(C_1\) is truth-free, so has a defined
classical value at \(K\). At any fixed point \(F\), leastness and
persistence give the same value; transparency makes its truth code have
that value. For \(d\in C_1\), \(T(d)\land A_1(d)\) therefore
has that defined value. For \(d\notin C_1\), the guard is false,
so the conjunction is false even when \(T(d)\) is undefined.
Thus \(T_1\) is bivalent on the entire domain, has the stated positive
extension, and is independent of \(F\).

Inductively suppose \(T_r\) is such a bivalent, fixed-point-independent
predicate. Every translated formula of \(L^{[r]}\) is built from
base atoms and atoms of this predicate by classical Boolean and quantifier
constructors. Induction on its formula tree shows that its values coincide
with the ordinary classical interpretation of \(L^{[r]}\), at every
fixed point, including \(K\). Hence every sentence code in
\(C_{r+1}\) has a defined fixed-point-independent value. Multiplying
by the next bivalent guard makes \(T_{r+1}\) false outside that set
and equal to its truth value inside it, exactly as in the first step.
This proves the induction and all the assertions. Being defined at
\(K\), each translated sentence is grounded. No uniform finite bound
on the ranks of all sentences at a level was used or follows from this
argument when its quantifiers range over infinitely many codes.
\end{proof}
\end{annexitem}

\begin{annexitem}{annex:p11.ktf.satisfaction}{thm:p11-ktf-satisfaction}{Least satisfaction fixed point}
\begin{proof}
Injectivity of both codings uniquely identifies the finite sequence and
formula of a valid pair. Coincidence of evaluation on free variables
makes the value independent of its arbitrary default entries. The same
formula induction as for truth evaluation establishes unique satisfaction
values and persistence when either coordinate grows. No formula obtains
two values. Valid pairs receiving a positive sign and negative sign are
therefore disjoint, and invalid pairs occur only negatively; \(\Psi\)
is a well-defined monotone operation on partial subsets of \(D^2\).

The recursion and increasing-state induction used for truth apply verbatim
with \(D^2\) in place of \(D\): at a limit, take the consistent
coordinate union; at a successor, use \(\Psi\). To prove a bound
without assuming that the argument domain is countable or well ordered,
encode a state by its signed subset
\((Q_1\times\{1\})\cup(Q_2\times\{0\})\) of
\(D^2\times\{1,0\}\). Suppose no stabilization occurred below
\(\eta=h(\mathcal P(D^2\times\{1,0\}))\). Any two stages
\(\alpha<\beta<\eta\) would be distinct: strict increase at
\(\alpha+1\leq\beta\) precludes equality with stage \(\alpha\).
Sending each ordinal \(\alpha<\eta\) to that encoded state would
inject \(\eta\) into the powerset defining its Hartogs bound, a
contradiction. Hartogs's theorem in ZF follows by collecting the order
types of all well-orderable subsets of the given set: these order types
form a set by Replacement, and the least ordinal exceeding every one
cannot inject into the set. This supplies the stated bound with no choice
of a new argument at each stage.

At stabilization the jump equality gives the displayed satisfaction
biconditionals. Induction shows each bottom stage below every fixed point;
the stabilized state is consequently least. Unlike truth on countably
many sentences, changes can involve the arguments of a free variable
throughout \(D\). Counting formula codes alone cannot bound their
number, so the earlier countable-code argument is unavailable for an
uncountable domain.
\end{proof}
\end{annexitem}

\begin{annexitem}{annex:p11.ktf.signed}{thm:p11-ktf-signed}{Signed closure equals the least truth fixed point}
\begin{proof}
Let \(V_K=\{(A,i):v_K(A)=i, i\in\{1,0\}\}\).
All required atomic signs belong to \(V_K\), including negative
non-code truth atoms at the fixed point. Every forcing condition for a
Boolean or quantified formula is its defining semantic clause. The
substitution lemma identifies closed numerical instances with the
corresponding assignments, and standard \(\mathbb N\) is exhausted
by those numerals. Transparency supplies the truth-prefix closure.
Thus \(V_K\) is a closed set, and leastness gives \(W\subseteq V_K\).

For the other inclusion induct on stages, proving that every defined
sentence value at \(K^\alpha\) has its sign in \(W\). At each
stage use a subordinate induction on formula complexity. At stage \(0\),
truth atoms are undefined; every defined ordinary atomic value is a
required base sign. Boolean defined values are forced by the cases
already specified. For a quantified value use the appropriate set of
closed numerical instances of its lower-complexity operand. For a
successor stage, if a truth atom has a non-code denotation, its negative
sign is a required sign. If it has a code denotation, that sign was
put into its coordinate by evaluation of the coded sentence at the
preceding stage, whose sign belongs to \(W\) by stage induction;
truth-prefix closure gives the sign of the atom. For a limit stage any
defined truth-atom membership came from an earlier stage and the same
reasoning applies. Ordinary atoms are unchanged at all stages.

For complete precision when a truth atom is \(T(t)\) with a closed
term rather than a numeral, its denotation is some natural number \(k\).
Equality in the base interpretation identifies its value with that of
\(T(\bar k)\). The specified atomic closure conditions include
replacement of equal-denoting closed terms; equivalently
one can first replace every closed argument by its numeral when applying
an atomic sign condition. This replacement concerns denotation, with
equality fixed and bivalent, and does not assert syntactic identity.
With this convention the truth-atom cases just given cover all closed
terms, including \(d,n,q\). Taking the stabilized stage gives
\(V_K\subseteq W\). Equality proves consistency of \(W\),
and transparency at \(K\) proves the converse truth-prefix closure.
\end{proof}
\end{annexitem}

\begin{annexitem}{annex:p11.ktf.schemes}{thm:p11-ktf-schemes}{Monotone alternatives and comparison}
\begin{proof}
For weak Kleene, a defined binary compound has all operands defined,
which stay defined with the same values under an information extension.
A defined quantified formula has every instance defined; their unchanged
values preserve the classical result. Negation and atoms satisfy the
same property. Formula induction gives persistence. For supervaluation,
the family of total extensions of \(R\) is a subset of the family
of total extensions of \(S\) when \(S\preceq R\). Both are
nonempty, since every partial state has a total extension. A unanimous
classical truth or falsity over the larger family persists over its
nonempty subfamily. Both jumps are consequently consistent and monotone.
The ordinal-recursion proof, including the countable-code selection,
applies to each without any change of its initial or non-code convention.

At a single state, every weak-Kleene defined value is also a Strong-Kleene
value of the same sign. This follows by structural induction, because
weak evaluation is defined only when all operands or instances are
defined, and both schemes then agree with classical evaluation.
Every Strong-Kleene defined value persists in every total extension by
the persistence lemma; there the evaluation is classical. It is therefore
the corresponding supervaluation value. Thus the three jumps at one state
satisfy \(\Phi_{\mathrm{wk}}(S)\preceq\Phi(S)\preceq
\Phi_{\mathrm{sv}}(S)\). Inducting on stages, combine these inequalities
with each jump's monotonicity to obtain the stagewise comparisons.
At limits take coordinate unions. At an ordinal beyond all three
countable closure stages this comparison is the asserted comparison of
the three least fixed points.

A classically valid sentence is true in every total extension of the
bottom state, hence supertrue there and positively assigned after one
jump. Classical soundness puts every \(\mathsf{PQT}\) code among
these positive codes at stage \(1\). In every total extension of that
stage, an instance with a true \(\mathsf{PQT}\) antecedent has a true
truth atom; an instance with a false antecedent has a classically true
implication. Its universal is true in all those completions, hence
supertrue at stage \(1\) and positively assigned at stage \(2\).
For the liar at a weak-Kleene fixed point the single negation
and atom have the same self-negation equation as under Strong Kleene.
For supervaluation, if its code is assigned true, every total extension
assigns that atom true, so the liar is superfalse, contradicting the
fixed point. If assigned false the reversed contradiction holds. If
unassigned, there are total extensions choosing its code positively
and negatively, so its liar evaluation takes both classical values
and is undefined. Therefore every supervaluation fixed point also
leaves its code and its own sentence value undefined.
\end{proof}
\end{annexitem}

\begin{annexitem}{annex:p11.ktf.restricted}{thm:p11-ktf-restricted}{Consistency-restricted fixed points}
\begin{proof}
For the negation-consistent family, eligibility is equivalent to
absence from \(S_1\) of a sentence code together with its negation's
code: necessity is immediate, and sufficiency is witnessed
by the total state with positive coordinate \(S_1\) and negative
coordinate its complement. For the maximal-consistent family,
eligibility is equivalent to the consistency of
\[
 \Sigma_S=\{A:c(A)\in S_1\}\cup
          \{\neg B:c(B)\in S_2\},\qquad S_1\subseteq C.
\]
Necessity uses that a complete consistent theory excluding \(B\)
includes \(\neg B\). For sufficiency, enumerate all sentences in
the countable language. Starting with \(\Sigma_S\), add each sentence
if that preserves consistency and otherwise add its negation. At least
one choice preserves consistency: if both additions yielded finite
contradiction proofs, the classical deduction and case principles would
yield a finite contradiction from the preceding set alone. A finite
contradiction from the union would occur already at one finite step.
Its deductive closure is consequently consistent and decides every
sentence. Interpreting \(T\) by its codes and making non-codes false
gives the required total completion. This is the countable Lindenbaum
argument; it uses the classical finite-proof notion of consistency.

At the bottom both eligibility conditions hold: the empty classical
theory is consistent, as witnessed by any classical structure, and the
preceding argument supplies a complete extension. When \(S\preceq R\)
are eligible, every admitted completion of \(R\) is an admitted
completion of \(S\). The unanimous defined values persist over this
nonempty subfamily. To see that a jump output is eligible, choose any
admitted total completion \(H\) of its input and evaluate the whole
language classically in the fixed base structure with \(T\) interpreted
by \(H\). Its classical truth set \(\operatorname{Th}(\mathcal M,H)\)
is a complete consistent theory, by classical soundness and bivalent
semantics. Its codes, with all non-codes made negative, define an admitted
total state of either kind. Every supertrue sentence at the input is
true in this chosen interpretation and every superfalse one is false
there. This total state thus extends the jump output, proving eligibility
and disjointness of its two coordinates.

If an increasing chain has eligible members, its union is still eligible.
For the maximal-consistent case every finite subset of its
\(\Sigma\)-set, including the negative-coordinate constraints, lies
in one chain member. A finite contradiction from that union would
therefore contradict the member's eligibility. The Lindenbaum argument
supplies a maximal-consistent completion of the union. For the
negation-consistent case, the two codes of an opposing sentence pair
would likewise occur together in one chain member, contradicting its
eligibility; a total completion is obtained by taking the union's
positive coordinate and its complement. The bottom iteration is
increasing by monotonicity, and
these observations give a valid state and eligibility at every limit.
The countable-code proof now gives stabilization, and comparison with
any eligible fixed point gives leastness. Thus these are fixed points
of their own restricted jumps, whose domains have been specified.

For any admitted negation-consistent completion, a sentence and its
negation cannot both belong to the positive coordinate. Thus every
sentence-code instance of \(\neg(T(x)\land T(n(x)))\) is true;
the false guard makes its implication true at non-sentence inputs.
This proves guarded universal supertruth at stage \(0\). After the
first jump every non-code belongs to the negative coordinate, so the
first truth atom is false for those arguments in all further extending
completions. The unguarded universal is therefore supertrue at stage
\(1\). For maximal-consistent completions the positive coordinate
contains only sentence codes from the outset, so its unguarded
universal is supertrue already at stage \(0\).
A maximal-consistent theory contains either \(A\) or
\(\neg A\) for every sentence \(A\), so each sentence-code
instance of \(T(x)\lor T(n(x))\) is true; non-sentence inputs are
guarded by the false \(\mathsf{Sent}\) antecedent. This proves its
universal supertruth. The jump definition gives the stated code-entry
offset in each case.

To check explicitly that no total fixed point has been smuggled into
this argument, a defined liar code at any eligible fixed point would
force the opposite sentence value in every extending completion, just
as in the unrestricted case. Either defined sign contradicts the jump
equality. Thus the liar remains undefined there, even though the
restricted maximal-consistent disjunction involving it is supertrue.
\end{proof}
\end{annexitem}

\begin{annexitem}{annex:p11.ktf.closed}{thm:p11-ktf-closed}{Conditional transparency after closing off}
\begin{proof}
Disjointness gives \(F_2\subseteq D\setminus F_1\), so
\(F\preceq\widehat F\). Persistence establishes that every
defined old value remains its value in the total extension, where
Strong-Kleene and classical evaluation coincide. At the fixed point
the code of \(A\) belongs to \(F_1\) exactly when \(A\) is
true, and the code of \(\neg A\) belongs to \(F_1\) exactly
when \(A\) is false. Thus the antecedent of the displayed conditional
is true at the closed-off state exactly when \(A\) was defined
at \(F\). In that case its retained classical value equals
membership of its code in \(F_1\), the value of \(T(k)\).
The biconditional is true. In the remaining case the antecedent is
false, so the classical conditional is again true.

The liar's code belongs to neither coordinate of \(F\), hence not
to \(F_1\). Closing off makes its truth atom false and its negation,
the liar, classically true. Its code still lies outside the positive
coordinate, so this total state fails the truth-jump fixed-point equation.
The set of classically true sentences of the closed-off language differs
from the old positive extension on this very code.
\end{proof}
\end{annexitem}

\begin{annexitem}{annex:p11.ktf.internal-syntax}{thm:p11-ktf-internal-syntax}{Elementary parsing and closing of formulas}
\begin{proof}
All numerical indices in this proof range over the specified standard
copy \(\mathsf N_{\mathcal A}\); all sequence objects range over
\(\mathsf{Seq}_{\mathcal A}\). Their successor relation is
elementary: \(m\) is the successor of \(n\) when \(n<m\)
and every \(i<m\) satisfies \(i\leq n\), with the order
restricted to this copy. For every fixed finite arity the tuple graph
is elementary. Its defining formula says that the output is a sequence
of that length with the specified coordinate values. Injection of
\(\iota\) gives uniqueness, including the zero-arity empty tuple.
References below to a tuple constructor mean this graph, not a new
primitive function.

First establish the finite-history recursion resource. If elementary
total functions \(g(\vec x)\) and \(h(y,k,\vec x)\) are given,
the function satisfying
\[
 f(0,\vec x)=g(\vec x),\qquad
 f(k+1,\vec x)=h(f(k,\vec x),k,\vec x)
\]
has an elementary graph on \(\mathsf N_{\mathcal A}\).
A witness is a coded sequence of length \(k+1\). Its first
entry is \(g(\vec x)\); for every numerical \(j<k\) its
entry \(j+2\) is obtained by applying the elementary graph of
\(h\) to entry \(j+1\), \(j\) and \(\vec x\).
Its last entry is the proposed output. These are one first-order
condition on the witness, using the elementary length and projection
graphs. Every standard finite recursion has this finite sequence,
and induction on the ordinary finite integer \(k\) proves
uniqueness. On an input outside the internal numerical copy define
the output to be \(0_{\mathcal A}\). This is another elementary
case. In particular addition, multiplication and the numerical
bookkeeping for finite certificates can be defined by such histories.
This is the actual recursion argument of 5A.1
\cite[pp.~65--66]{Moschovakis1974EIASOriginal}.

A term code has one of the two displayed fixed tuple forms. Its
elementary denotation relation \(\mathsf{Den}(t,s,a)\) says
either that \(t\) is the variable code for some numerical
\(i\geq1\) and \(a=q_{\mathcal A}(s,i)\), or that it
is the constant code with payload \(a\). The two cases are
disjoint and their values unique. For valid assignment codes they
are exactly the declared term semantics, including default values.

Here is a finite certificate for being a formula. A coded nonempty
sequence of records is in postorder. Each record stores its proposed
formula code and, where appropriate, one or two indices of earlier
records. Its first coordinate must be a prime code, or the specified
negation, conjunction or existential constructor applied to the codes
at those earlier records. In a prime record the term codes have the
two valid forms just described; in an existential record the binder
index is a positive internal integer. The root code is the code in
the last record. The finitely many ordinary predicate cases are
explicitly selected by \(j=1,\ldots,l\), with their specified
arities, and the equality and \(T\) cases are separate. A single
first-order universal over the record indices verifies these finite
cases and their projections. Thus existence of such a certificate
is an elementary predicate.

Every finite formula supplies this certificate by traversing its
finite parse tree. Conversely induction on record indices constructs
an actual finite expression at every record. Because tuple coding
is injective and the constructor tags are distinct, its terms and
constructor children are unique, hence the root expression with that
code is unique. Domain objects which satisfy a cyclic constructor
equation do not pass a finite postorder certificate merely by satisfying
that equation: a child is required to occur at a strictly earlier
record. A constant payload is never treated as a child to be parsed.
This proves both the formula predicate and the injectivity asserted
in the language convention.

For the sentence predicate add to each record a finite increasing
sequence of positive numerical indices, listing its free variables.
The elementary verification is explicit. A variable contributes its
singleton index and a constant the empty list. A prime record takes
the union of its terms' lists; negation copies the list; conjunction
takes the union of the two child lists; an existential record deletes
its binder index. The union and deletion checks say, for every
internal numerical index, that membership in the proposed finite
list is respectively membership in either input list or membership
in the input list unequal to the deleted index. Increasing order
makes the list unique. These conditions are first order using the
projection and order graphs. Finite induction shows that the root
list is exactly the actual free-variable list. It is empty exactly
for a sentence. The same construction gives the elementary graph
of the greatest free index, taking zero for an empty list.

Assignment update at positive index \(i\) to value \(a\)
has the following elementary graph. Its output is a sequence of
length \(\max(\ell_{\mathcal A}(s),i)\). At position \(i\)
its value is \(a\); at every other position \(j\) its value
is \(q_{\mathcal A}(s,j)\). Existence follows by taking
this finite tuple; injection gives uniqueness. Its assignment semantics
is exactly replacement of variable \(v_i\), retaining defaults
outside the new finite range.

To substitute \(\dot a\) for free \(v_i\), annotate a
postorder parse certificate with an output code for each record.
The variable term with index \(i\) is replaced by the constant
term code with payload \(a\); other terms are copied. Prime
records rebuild their tuple with these output terms; negation and
conjunction rebuild with their children's output codes. An existential
with binder \(i\) copies its entire original code unchanged;
an existential with a different binder rebuilds using the child's
output code. Each rule is a finite elementary case. Finite induction
proves existence and uniqueness of the output and identifies it
with free substitution. Copying the whole matching-binder expression
is essential: occurrences bound there are not replaced. A substituted
constant has no variable occurrences to be captured.

Let \(m\) be the greatest free index of the input formula.
Starting with its code, perform the preceding free substitution for
\(i=1,\ldots,m\), using the constant with value
\(q_{\mathcal A}(s,i)\). The elementary finite-history recursion
already proved verifies this sequence of output codes. Every free
variable has been replaced; no new free variable has been introduced.
The final output is therefore a sentence code, and its interpretation
is the original formula under the assignment denoted by \(s\).
For invalid formula or assignment inputs choose the fixed code of
\(\dot 0_{\mathcal A}\ne\dot 0_{\mathcal A}\).
With this additional elementary case the graph is total and unique.
This supplies \(\mathsf{Cl}_{\mathcal A}\) and completes all
the assertions without assuming a uniform first-order truth predicate.
\end{proof}
\end{annexitem}

\begin{annexitem}{annex:p11.ktf.internal-induction}{thm:p11-ktf-internal-induction}{The interpreter is one elementary positive induction}
\begin{proof}
The syntax lemma supplies all nonrecursive guards and code graphs.
For clarity spell out the potentially nonfinite rule. If \(e\)
codes \(\exists v_i B\), its true rule at \(s\) is
\[
 \exists a\exists s'\,
 (\mathsf{Upd}(s,i,a,s')\land
       \mathsf V(c(B),s',1_{\mathcal A})),
\]
and its false rule is
\[
 \forall a\exists s'\,
 (\mathsf{Upd}(s,i,a,s')\land
       \mathsf V(c(B),s',0_{\mathcal A})).
\]
The update graph is total and unique, so these conditions are the
required witness and all-instances conditions. No negative occurrence
of \(\mathsf V\) is introduced. A truth atom's rule uses
\(\exists a(\mathsf{Den}(t,s,a)\land\mathsf Z(a,b))\).
For an ordinary atom the denotation graph supplies its finitely many
arguments and the corresponding fixed \(R_j\) or its complement,
or equality or inequality, supplies the sign. Ordinary negations here
are permitted in an elementary induction. A conjunction uses the
two positive true premises or the disjunction of its two false premises.
Negation interchanges the two marker constants. The sentence and
nonsentence \(\mathsf Z\) rules have the elementary guards
\(\mathsf{Sent}_{\mathcal A}\) and its complement.

Pack the two relations into a single four-ary relation \(P\).
Use \(P(0_{\mathcal A},e,s,b)\) for \(\mathsf V(e,s,b)\)
and \(P(1_{\mathcal A},u,\iota(),b)\) for
\(\mathsf Z(u,b)\). Take the union of the translated rules,
with their indicated tags and valid-code guards, and the retention
rule. This is one positive first-order formula in \(P\).
All untyped tuples remain outside its least fixed point. Restricting
that point to the two tags recovers exactly the least simultaneous
rule closure, say \(W_V,W_Z\). Packing and unpacking the rules
proves this equality in both directions, or directly by the leastness
of the two closures. Parameter sections implement the fixed tags.

Let \(F=K_{\mathcal A}^{\mathrm K}\), whose ordinal existence
uses information monotonicity and the stated Hartogs bound. The two
relations consisting of all actual formula signs evaluated at \(F\)
and its truth-predicate signs satisfy every displayed rule. Formula
semantics verifies the evaluation rules; the truth fixed-point equation
verifies sentence copying and the nonsentence false rule. Therefore
leastness gives \(W_V,W_Z\) as subsets of these semantic relations.
In particular \(W_Z\) is consistent, and \(W_V\) has no
opposed signs for a given formula and assignment.

Let \(S\) be the partial truth state specified by \(W_Z\).
Induction on the actual finite formula shows
\[
 (e,s,b)\in W_V
 \quad\Longleftrightarrow\quad
 \text{the formula coded by }e\text{ has sign }b\text{ at }S.
\]
For the reverse implication the requisite premises already belong to
\(W_V,W_Z\) by the induction hypothesis and the closure rules
add the conclusion; the universal false-existential rule uses all
its instances, not a finite list of them. For the forward implication
consider the least rule-generation stages. A generated formula sign
has one of the displayed semantic justifications; retention cannot
originate a sign, since generation starts empty. Once true of \(S\)
these justifications remain true. Equivalently induction on the formula
works together with its sign's original generating rule. This proves
the assertion for every assignment. Existence of such a generation
sequence is the ordinary bottom iteration of the monotone rule operator;
its length is bounded by the Hartogs ordinal of its ambient powerset.

Sentence copying and the nonsentence rule now give
\(J_{\mathrm K}(S)\preceq S\). Conversely every
\(W_Z\)-sign originates in one of those two rules, whose
evaluation premise is a sign at \(S\). Thus
\(S\preceq J_{\mathrm K}(S)\) and \(S\) is a fixed
point. Leastness of \(F\) gives \(F\preceq S\), while the
earlier semantic containment gives \(S\preceq F\). Hence \(S=F\),
and the formula assertion identifies \(W_V\) exactly as claimed.

Finally fix a formula \(B(\vec v)\) and its finite domain
parameters. Add to the finite positive system an output component
whose only rule copies
\(\mathsf V(c(B),\iota(\vec a),b)\) to the tuple
\(\vec a\). The fixed-arity tuple graph is elementary, so
this is a positive elementary rule. It does not feed back into the
other components; its least closure is exactly the desired extension.
The same tagged padding construction packs these finitely many
components into one relation. Thus both sign extensions are inductive,
including elementary preimages and additional parameter sections,
without assuming a separate closure or representation theorem.
\end{proof}
\end{annexitem}

\begin{annexitem}{annex:p11.ktf.acceptable-ind}{thm:p11-ktf-acceptable-ind}{Strong-Kleene definability over an acceptable structure}
\begin{proof}
The interpreter lemma proves containment of every extension in
\(\operatorname{Ind}(\mathcal A)\). For the reverse containment
let \(\Gamma\) be induced by a positive formula
\(\phi(\vec x,S)\) with its finitely many parameters. Move
negations to ordinary atoms; since \(S\) occurs positively,
none of its occurrences is negated. The other connectives and
quantifiers are the positive operations \(\land,\lor,\forall,\exists\).

For a prospective template code \(k\) and tuple \(\vec z\)
define the elementary graph
\[
 G(k,\vec z,u)\quad\Longleftrightarrow\quad
 u=\mathsf{Cl}_{\mathcal A}(k,\iota(k,\vec z)).
\]
This notation expands to the already proved closing and tuple graphs.
Form a template \(B(v_1,\vec x)\) by replacing every
\(S(\vec t)\) in \(\phi\) by
\(\exists u(G(v_1,\vec t,u)\land T(u))\), using fresh bound
variables throughout. Conjoin \(v_1=v_1\) and
\(\bigwedge_j x_j=x_j\). These redundant ordinary conjuncts
ensure that its designated argument variables actually occur freely,
even for a constant operator or an unused coordinate. The graph
formula is in the base language and has finitely many parameters.
Let \(b=c(B)\) and define the actual sentence
\(\theta_{\vec a}=B(\dot b,\dot{\vec a})\).
Because every object has a constant name, this is an object-language
sentence. Its code is \(\mathsf{Cl}_{\mathcal A}
(b,\iota(b,\vec a))\). No graph has been assumed to define
\(\Gamma\)'s desired fixed point: it performs only finite syntax
closing of its prospective template argument.

In \(\theta_{\vec a}\) the graph at an \(S\)-occurrence
names exactly the code of \(\theta_{\vec t}\), since the
template argument has become the fixed name \(\dot b\).
The graph has one output and is bivalent. Therefore the replacing
formula is true exactly when that sentence code belongs to the
current positive truth coordinate. Positive-formula induction gives,
for every truth state \(Q\),
\[
 v_Q(\theta_{\vec a})=1
 \quad\Longleftrightarrow\quad
 \vec a\in\Gamma(\{\vec z:c(\theta_{\vec z})\in Q_1\}).
\]
The existential graph does not create extra true witnesses: only
its unique output can have a true guard. An ordinary literal is
bivalent, and the true conditions of each positive connective or
quantifier are exactly its corresponding classical condition on the
positive extension. Undefined truth atoms count as nontrue in these
conditions, just as required by the set on the right.

Consequently, if \(Q^\alpha\) is the original bottom truth
iteration and \(E^\alpha=\{\vec a:c(\theta_{\vec a})
\in Q^\alpha_1\}\), then
\(E^0=\varnothing\), \(E^{\alpha+1}=\Gamma(E^\alpha)\),
and \(E^\lambda=\bigcup_{\alpha<\lambda}E^\alpha\)
at a limit. The latter equality uses membership of a fixed code
in a coordinate union. It does not require countability. Transfinite
induction identifies this with \(\Gamma\)'s bottom iteration,
so its eventual value is its least fixed point. The formula
\[
 \exists u\,
 (G(\dot b,\vec x,u)\land T(u))
\]
has exactly that positive extension at
\(K_{\mathcal A}^{\mathrm K}\). Fixing further coordinates
at their named parameter values gives every allowed parameter section.
The same construction works in every finite arity.

For a formula defined at every input, its negative extension is
the complement of its positive extension. Both were proved inductive
by the signed interpreter. Their intersection-of-classes condition
is exactly the definition of hyperelementary. This proves the final
implication without claiming its as yet separate totality converse.
\end{proof}
\end{annexitem}

\begin{annexitem}{annex:p11.ktf.elementary-games}{thm:p11-ktf-elementary-games}{From induction to games and from games to truth}
\begin{proof}
Let a positive elementary formula \(\phi(\vec x,S)\)
induce \(\Gamma\). Put its finite expression in negation normal
form, with negations only on ordinary atoms. Use one control label
per subformula, storing the current values of all the finitely many
variables occurring in \(\phi\). Registers include its input
coordinates and all its bound-variable coordinates. Ordinary literals
are terminal tests: accept when the bivalent literal holds and reject
otherwise. A conjunction moves to a binary minimum node with the
two subformulas and the same register tuple; a disjunction moves to
a binary maximum node. An existential or universal moves to a
domain maximum or minimum node whose successor updates that
bound-variable register and selects the matrix label. An
\(S(\vec t)\) node makes a deterministic recursive call to
the root label, setting its input registers to the term values and
resetting the other registers to the specified default. No return
stack is required: the recursive predicate's requested value is
the value of that called configuration itself. All tests and register
updates are elementary finite cases; the relational language's terms
are variables or constants. Fixed-arity sequence coding makes the
valid configuration set and each successor graph elementary.

Replace each rejecting terminal by a deterministic self-loop.
Call this resulting game \(P\). Its least labeling contains
only \(1\) and \(\mathbf u\): induction on stages shows
that no rule introduces \(0\), while a self-loop starting
undefined remains undefined. Accepted configurations are consequently
the least positive rule closure of the true rules for the same finite
control diagram. Those rules at an ordinary test accept exactly
its true instances; at a conjunction they require both children,
at a disjunction either, at an existential some domain choice, and
at a universal all choices. At a recursive-call node they copy
the root acceptance with its designated tuple. Replacing rejection
by divergence changes none of these true rules.

To identify the root relation let \(Y\) be this least rule
closure and let \(B\) be its root-input relation. Induction
on the fixed finite subformula tree, using the recursive-call rule
as its atomic \(S\) case, shows that each subformula label
is accepted precisely when that subformula is classically true with
\(S=B\) and the current registers. In the forward direction
every accepted label originates in one of its indicated true rules;
in the reverse direction closure applies the required rule after
the inductively identified child premises. Thus \(B=\Gamma(B)\).
If \(I\) is any fixed point of \(\Gamma\), assign to every
subformula configuration its classical true condition with \(S=I\).
These sets form a closed positive control system, including the
recursive call at the root. Leastness gives \(B\subseteq I\).
Therefore \(B\) is the least fixed point of \(\Gamma\).
Named fixed coordinates implement parameter sections. This proves
the first assertion by an explicit compilation, without importing
a programme characterization of induction as an assumption.

For the second assertion let any finite elementary-control game
be given, encoding its valid states by objects \(z\in A\).
Set
\(G(k,z,u)\) to be the elementary graph of
\(u=\mathsf{Cl}_{\mathcal A}(k,\iota(k,z))\).
Build a template \(B(v_1,v_2)\) with a finite ordinary-guarded
disjunction selecting the type of configuration \(v_2\).
The accept case is true and the reject case false. A deterministic
case takes the unique elementary successor \(w\) and uses
\(\exists u(G(v_1,w,u)\land T(u))\). A binary maximum
is the disjunction of that expression for its two successors;
a binary minimum is their conjunction. A domain maximum or
minimum quantifies \(\exists a\) or \(\forall a\)
over the successor graph and that expression. Each successor graph
is total and unique on its guarded case. An existentially quantified
unique graph therefore has exactly the value of the requested
truth atom, even when that value is undefined. A universally
chosen domain successor has the minimum of those values. False
guards in the other cases contribute false to the outer disjunction
and do not disturb the selected value. Conjoin \(v_1=v_1\)
and \(v_2=v_2\) to retain both template coordinates.

Let \(b=c(B)\), and for a valid configuration \(z\) put
\(\theta_z=B(\dot b,\dot z)\). Its code is the stated
closing operation, and the graph inside it calls precisely
the code of its child's \(\theta_w\). At stage zero all
these codes have undefined truth-predicate values. The truth jump's
values on these codes at stage \(\alpha+1\) are precisely
the game's update of its values at stage \(\alpha\).
This includes terminal codes, which acquire their signs at jump
one, and self-loops, which remain undefined. At a limit both
constructions use coordinate information unions. Transfinite induction
therefore identifies the game labels at every stage with these
truth-code labels. It does not use the truth values of the terminal
sentences at stage zero as if their codes were already defined.
At the least truth fixed point the displayed truth atoms have
exactly the game's least labels.

For an input tuple \(\vec a\), its initial configuration
\(z(\vec a)\) has an elementary total graph given by its
fixed control label, the input registers and fixed default registers.
The formula using that graph followed by
\(\exists u(G(\dot b,z,u)\land T(u))\) realizes the
desired label at each input. All added parameters are the finitely
many coding and game parameters; names for them are available
in the stated language. No restriction on \(|A|\) or a
uniform interpreter for an infinite primitive signature was used.
\end{proof}
\end{annexitem}

\begin{annexitem}{annex:p11.ktf.selective-shuffle}{thm:p11-ktf-selective-shuffle}{A totality-producing shuffle}
\begin{proof}
The shuffled state is \((\varepsilon,c,d)\), where
\(\varepsilon\in\{0_{\mathcal A},1_{\mathcal A}\}\)
selects \(P\) or \(Q\), respectively. If the active state
is terminal, return its sign. Otherwise take its deterministic,
binary or domain-choice transition, retain the other component,
and switch \(\varepsilon\). Its maximum or minimum type
is retained for that one transition. This is a finite elementary
partition of a fixed-arity encoded configuration set, with elementary
total successor graphs. The shuffle thus belongs to the stated
class and has its own least information labeling.

Here is a rank induction proving the required sign behavior. For
each proper resolution rank \(r\), establish both claims:

\begin{enumerate}
\item An active component resolved at rank \(r\) wins with
its sign if the other component's rank is at least \(r\),
including the unresolved marker.
\item An inactive component resolved at rank \(r\) wins with
its sign if the active component's rank is strictly greater than
\(r\), including the unresolved marker.
\end{enumerate}

First prove claim one, using both claims at smaller ranks. A terminal
active component returns its sign. A deterministic winning component
has a child of the same sign with rank strictly below \(r\).
After its move that child is inactive and the competitor has rank
at least \(r\); claim two at the smaller rank applies. For
a \(P\) maximum with label \(1\), at least one child
has label \(1\) at a rank below \(r\), so its shuffled
maximum is \(1\). For a \(P\) minimum with label \(1\),
all children have label \(1\) below \(r\), so all the
shuffled children are \(1\) and their minimum is \(1\).
The \(Q\) cases are dual: a maximum labeled \(0\) has
all children labeled \(0\) below \(r\); a minimum labeled
\(0\) has at least one such child. Claim two at those smaller
ranks supplies the shuffled signs. These arguments cover binary
and domain choices. If all the required shuffled children resolve,
their ranks have an ordinal supremum because they form a set,
so the universal/minimum or universal/maximum sign is indeed generated
at some stage of the least labeling.

Next prove claim two at the same \(r\), using claim one
just established. Write \(r=\delta+1\), which is possible
because proper resolution ranks are successors. Suppose the winner
is inactive \(P\) with sign \(1\), and active \(Q\)
has rank greater than \(r\). A deterministic \(Q\) child
cannot have a rejecting rank below \(r\), since that would
reject its parent by stage \(r\). Thus the child has rank
at least \(r\) or is unresolved. After the move, claim one
at \(r\) makes active \(P\) win.

At a \(Q\) minimum, every child's rejecting rank must likewise
be at least \(r\), or undefined: a smaller rejecting child
would already force the parent to reject by stage \(r\).
So \(P\) wins after every such transition and the shuffled
minimum is \(1\). At a \(Q\) maximum, at least one
child has rank at least \(r\) or is undefined. Otherwise
all children would already reject by stage \(\delta\), and
the parent would reject by \(\delta+1=r\), a contradiction.
The corresponding shuffled branch is \(1\), so the maximum
is \(1\). An active rejecting terminal has rank one and
cannot have rank greater than \(r\geq1\); an accepting
\(Q\) terminal is excluded by the hypothesis. This exhausts
the \(Q\) cases.

If the inactive winner is \(Q\) with sign \(0\), the
active competitor is \(P\). At a \(P\) maximum all
children have accepting rank at least \(r\) or are unresolved,
since an earlier accepting child would resolve that maximum by
\(r\). After every transition claim one makes \(Q\)
win, so the shuffled maximum is \(0\). At a \(P\)
minimum at least one child has rank at least \(r\) or is
unresolved; otherwise all would accept by \(\delta\) and
the parent by \(r\). That branch has shuffled value \(0\),
so the minimum is \(0\). Deterministic and terminal cases
are the same dual cases already considered. This proves claim two.

Transfinite induction now proves both claims at every proper rank.
With exactly one initial resolved component, its rank is smaller
than the other's unresolved marker. Claim one or two applies
according to which component is initially active, yielding the
specified sign. This is an eventual-resolution proof. It does not
assert a running-time equality obtained by confusing
\(\sup_i r_i+1\) with \(\sup_i(r_i+1)\).
The successor form of \(r\) is precisely what justified the
all-children bounds in the proof.
\end{proof}
\end{annexitem}

\begin{annexitem}{annex:p11.ktf.hyperelementary}{thm:p11-ktf-hyperelementary}{Total Strong-Kleene definitions are exactly hyperelementary}
\begin{proof}
The acceptable-induction theorem already proves that a total formula's
positive extension and its complement are inductive, giving the
upper inclusion. Conversely let \(R\) be hyperelementary.
Apply the explicit induction-to-game construction to \(R\)
and its complement, including their permitted parameter sections.
Obtain games \(P\) and \(P'\) which only accept or remain
undefined, with accepted initial inputs exactly \(R\) and
its complement.

Dualize \(P'\) by interchanging accept and reject, maximum
and minimum, and retaining its deterministic transitions and tests.
The least labels are exchanged \(1\leftrightarrow0\) and
\(\mathbf u\leftrightarrow\mathbf u\), as induction on
stages proves directly. Thus this dual game \(Q\) only
rejects or remains undefined, and rejects exactly the complementary
inputs. At every input exactly one of the initial configurations
of \(P,Q\) resolves: the \(P\) one on \(R\), and the
\(Q\) one off \(R\). Their encoded input registers may
be kept disjoint in the shuffle; each game retains its own defaults
and finite parameters.

The selective shuffle lemma makes the initial shuffled label \(1\)
on \(R\) and \(0\) off \(R\). The game-to-truth
construction realizes these actual labels by a formula at
\(K_{\mathcal A}^{\mathrm K}\). Its positive extension
is \(R\) and it is defined at every input. Fixed-arity encoding
and the same input construction handle any finite arity. The argument
requires the genuine alternating shuffle, not an assumed total
combination of two gap-valued formulas.
\end{proof}
\end{annexitem}

\begin{annexitem}{annex:p11.ktf.general-classical-valuation}{thm:p11-ktf-general-classical-valuation}{Uniform classical evaluation with a relation argument}
\begin{proof}
Take the finite ordinary-predicate cases of the elementary interpreter,
using its denotation and parse graphs. Require \(V\) to have
exactly one of the two marker signs for every formula code and
sequence assignment, and to have no other tuples. In the \(T\)
case replace the recursive sign predicate by membership of the
unique term value in \(R\). Ordinary atoms use their bivalent
predicates. For a negation require the opposite sign of its child;
for a conjunction require the ordinary Boolean condition on its
children; for an existential require that the true sign holds exactly
when there exist a domain value and its uniquely updated assignment
with a true matrix sign. Totality then gives its false sign in the
complementary case. Quantifying over the relevant codes, assignments
and domain values gives one finite first-order matrix with only
\(R,V\) as free relation variables. The fixed signature is
finite, so this is not an infinite collection of primitive cases
hidden in one formula.

Existence is classical evaluation of the actual finite formulas under
every finite assignment and \(T=R\). Its graph is a set of
triples in \(A^3\): the set of formula codes is a subset of
\(A\), and assignment codes are another such subset. Recursion
on natural formula complexity defines these values, and the resulting
relation satisfies all the equations. For uniqueness induct on the
actual finite expression. Atomic equations force its sign; the
constructor equations force each further sign from the uniquely
identified child signs under every assignment. Every formula has
a finite parse certificate, so this covers all the tuples allowed
by the matrix. It proves the asserted uniqueness without identifying
this relation-argument matrix with an elementary absolute satisfaction
predicate. The classical positive-induction construction in source
5B.2 is compatible with this valuation relation, with its opposite
source flag convention already distinguished in the body.
\end{proof}
\end{annexitem}

\begin{annexitem}{annex:p11.ktf.general-tellers}{thm:p11-ktf-general-tellers}{Free truth-teller coordinates at the least supervaluation point}
\begin{proof}
Closing \(B\) at \((b,a)\) replaces its designated free
variables by \(\dot b,\dot a\), so its code is
\(\mathsf{Cl}_{\mathcal A}(b,\iota(b,a))\).
The graph \(G\) inside that very instance therefore selects
this same output. The ordinary equality conjuncts are true; the
unique graph guard is true exactly at its output. In any classical
interpretation \(T=R\), the value of \(\tau_a\) is thus
membership of its own code in \(R\). Retaining the second
equality conjunct ensures that the constant payload \(a\)
actually appears in the sentence. Injectivity of the finite-expression
coding therefore makes \(a\mapsto c(\tau_a)\) injective.
Its graph is the proved elementary closing graph.

Suppose all these coordinates are undefined at a stage
\(Q=(Q_1,Q_2)\). For any chosen subset \(B_0\subseteq A\)
put
\[
 R_{B_0}=Q_1\cup\{c(\tau_a):a\in B_0\}.
\]
The added codes belong to neither \(Q_1\) nor \(Q_2\).
Thus \(R_{B_0}\) and its complement form a total extension
of both coordinates of \(Q\); injectivity gives exactly the
selected truth-teller assignments. For each fixed \(a\), choose
\(B_0\) containing or excluding \(a\). Its sentence then
has both possible classical values in the permitted extension family.
Neither sign is supervaluationally forced, so its code stays undefined
at the next stage. At stage zero every truth-predicate coordinate
is undefined, and at a limit an absent coordinate remains absent
from each union. Transfinite induction proves the claim at every
stage and hence at the least point. The same explicit set construction
at that point proves the free-subset assertion. This direction uses
no selection of a family of arbitrary witnesses and no maximal-theory
or consistency restriction on completions.
\end{proof}
\end{annexitem}

\begin{annexitem}{annex:p11.ktf.general-sv}{thm:p11-ktf-general-sv}{General supervaluation definability}
\begin{proof}
For the upper direction work first in the complete powerset of
typed signed pairs \(A\times\{0_{\mathcal A},1_{\mathcal A}\}\).
Its elements \(X\) may be inconsistent; this ambient complete
lattice is different from the original consistent information poset.
Let \(\mathsf{Ext}(X,R)\) require, for every \(a\), that
\(X(a,1_{\mathcal A})\) implies \(R(a)\) and
\(X(a,0_{\mathcal A})\) implies \(\neg R(a)\).
Define an ambient operator \(H\) whose sentence-code sign
\((e,b)\) holds exactly when
\[
 \forall R\forall V\,
 ((\mathsf{Ext}(X,R)\land
       \mathsf{CV}_{\mathcal A}(R,V))
       \longrightarrow V(e,\iota(),b)).
\]
The nonsentence false signs are added as base outputs; no nonsentence
true sign is output. This is monotone because enlarging \(X\)
reduces its total extensions. Its bottom iteration has consistent
values: a consistent \(X\) has the explicit total extension
\(R=\{a:X(a,1_{\mathcal A})\}\), so no sentence can
be forced to both signs, and the nonsentence cases are disjoint
from sentence codes. The ambient operator on a consistent input
is exactly the original unrestricted supervaluation jump. Hence
its least fixed point is the signed version of
\(K_{\mathcal A}^{\mathrm V}\). Existence and leastness
follow by the ordinary monotone powerset iteration; no complete-lattice
theorem is being applied to the smaller consistent-state order.
Equivalently its least point is the intersection of all typed
pre-fixed sets \(H(X)\subseteq X\): iteration from the
empty set lies in every such set, and the eventual fixed point
is itself one of them.

The pre-fixed condition has a first-order witness matrix after
packing countermodels. For every sentence-code sign \((e,b)\)
missing from \(X\), \(H(X)\subseteq X\) requires a
total extension \(R_{e,b}\) and its valuation \(V_{e,b}\)
with \(\neg V_{e,b}(e,\iota(),b)\). Conversely these
countermodels, together with all nonsentence false signs, imply
the pre-fixed condition. Choice selects the family of witnesses
for this set of missing pairs. Pack it into one six-ary relation
\(Y\), using
\[
 \begin{split}
 R_{e,b}(a)&\Longleftrightarrow
 Y(0_{\mathcal A},e,b,a,0_{\mathcal A},0_{\mathcal A}),\\
 V_{e,b}(f,s,d)&\Longleftrightarrow
 Y(1_{\mathcal A},e,b,f,s,d).
 \end{split}
\]
Here the tags are distinct and the padding is fixed. Let
\(\mathsf{Pre}_{\mathcal A}(X,Y)\) be the finite first-order
matrix saying that \(X\) is typed, contains all nonsentence
false signs, and that for every missing sentence sign these
sections satisfy \(\mathsf{Ext}\), \(\mathsf{CV}_{\mathcal A}\)
and the displayed countervalue. Then
\[
 H(X)\subseteq X
 \quad\Longleftrightarrow\quad
 \exists Y\,\mathsf{Pre}_{\mathcal A}(X,Y).
\]
The forward implication is exactly the identified choice use;
the reverse implication uses the supplied countermodel sections.
No numerical canonical-choice argument is used on \(A\).

Now fix a formula \(F(\vec v)\) with its finite parameters,
and let \(e=c(F)\). Its supertrue extension at the least
point is defined by
\[
 \forall X\forall Y\forall R\forall V\,
 \bigl((\mathsf{Pre}_{\mathcal A}(X,Y)
       \land\mathsf{Ext}(X,R)
       \land\mathsf{CV}_{\mathcal A}(R,V))
       \longrightarrow V(e,\iota(\vec a),1_{\mathcal A})\bigr).
\]
The fixed-arity assignment tuple is expressed by its elementary
graph if it is not written as a term. This is a universal relation
block with a first-order matrix and finitely many individual
parameters. If \(F\) is supertrue at the least point, every
pre-fixed \(X\) contains that point, so every total extension
of \(X\) is a permitted extension of the least point and
the displayed implication holds. Conversely take \(X\) equal
to the least point and a witness \(Y\) to its pre-fixedness.
The implication then requires the true value in all its extending
classical valuations, exactly supertruth. This proves the upper
inclusion directly for formula extensions, without a negated
substitution of the least-point definition into an extension matrix.

For the lower direction write a given definition as
\(\forall U_1\cdots\forall U_m\,\psi(\vec a,\vec U)\),
where \(\psi\) is elementary and each \(U_j\) has finite
arity \(n_j\). Use the disjoint injective payload maps
\[
 p_j(\vec z)=\iota(j_{\mathcal A},z_1,\ldots,z_{n_j}).
\]
Their fixed-arity graphs are elementary, including the zero-arity
case. Replace each atom \(U_j(\vec t)\) in \(\psi\)
by the elementary graph taking \(\vec t\) to \(p_j(\vec t)\)
and then to \(c(\tau_{p_j(\vec t)})\), followed by
the truth atom at that code. Use fresh bound variables and retain
all ordinary negations and quantifiers. Call the resulting fixed
truth-language formula \(\widehat\psi(\vec v)\).
In a total extension, its replacing atom is precisely the Boolean
value assigned to that truth-teller code; graph uniqueness eliminates
all other witnesses.

The truth-teller lemma makes every subset of their payload domain
available among total extensions of the least point. In particular,
given arbitrary \(U_1,\ldots,U_m\), select the explicit subset
\(B_0=\bigcup_j p_j[U_j]\) and use its explicit extension
from that lemma. Disjoint tags and tuple injection reproduce all
the selected relations. Conversely every total extension induces
some such tuple of relations by its assigned values on these
codes. Finite classical formula induction identifies the value
of \(\widehat\psi(\vec a)\) with \(\psi(\vec a,\vec U)\)
for the induced relations. Thus its value is true in every permitted
extension exactly when the original universal-relation definition
holds. Its supertrue extension is exactly the requested relation.
This constructs every \(\Pi^1_1(\mathcal A)\) relation,
with its finite individual parameters and in every finite arity.
The lower construction itself needs no choice of arbitrary
countermodels; the theorem's choice premise is used only in the
upper packing argument.
\end{proof}
\end{annexitem}

\begin{annexitem}{annex:p11.ktf.general-finite-hierarchy}{thm:p11-ktf-general-finite-hierarchy}{Elementary finite Tarski comparisons over an acceptable structure}
\begin{proof}
For each fixed finite auxiliary vocabulary extend the internal
formula constructors by distinct tags for its finitely many truth
tokens, with an internal numerical level index and a term code.
Use the same term constructors and ordinary formula constructors.
The postorder certificate from the syntax lemma, with these finitely
many added prime cases, gives its elementary sentence predicate.
Its variable-term code \(c(v_1)=\iota(0_{\mathcal A},1_{\mathcal A})\)
is not a formula code: every prime or compound formula has a
different constructor tag. This holds for all the auxiliary vocabularies.
No claim that the object \(0_{\mathcal A}\) is outside
their code sets is required.

A translation which replaces finitely many specified atom patterns
by fixed formulas has an elementary code graph. To see this directly,
annotate a postorder input parse certificate with its proposed output
codes. Ordinary primes copy their codes; constructors rebuild with
their children's output codes. At a substituted prime the fixed
replacement's argument variable is replaced by its input term,
and its newly bound variables are chosen above the greatest variable
index in the input and replacement. The finite free-index computations
and numerical finite histories of the syntax lemma verify this
fresh renaming and term substitution. Existence and uniqueness of
the annotated outputs follow by finite parse induction, and all
verification conditions are elementary. Sending invalid sentence
inputs to \(c(v_1)\) makes the graph total. This argument
allows opaque domain constants in the input and keeps every fixed
replacement's finite parameters; it does not turn all constants
into numerical word symbols.

Now construct the finite guard sequence successively. The truth-free
sentence predicate defines \(C_1\) and hence a formula \(A_1\).
Once the earlier guard formula \(A_r\) is fixed, replace
each \(U_r(t)\) prime by \(T(t)\land A_r(t)\).
The preceding annotated certificate proves that this translation
has an elementary graph. Its image on the elementary auxiliary
sentence set is elementary, by first-order existential quantification
over its input and output graph. This defines \(C_{r+1}\)
and permits a fixed truth-free defining formula \(A_{r+1}\).
Induction on the finite level therefore supplies all the actual
guards demanded by Definition~\ref{def:p11-ktf-tarski-guards}.
The proof of Theorem~\ref{thm:p11-ktf-tarski} then applies to
the same guarded formulas and gives bivalence and fixed-point-independent
truth, since it uses no countable-code bound.

For completeness give the cumulative/noncumulative translations at
this nonnumerical domain. At a fixed level let \(\mathsf N_i\)
be the truth predicate for the preceding noncumulative language,
and \(\mathsf C_i\) that for the preceding cumulative language.
Each is classically false on other inputs. The code map
\(a\mapsto c(\mathsf N_i(\dot a))\) has an elementary
graph by the fixed constructors. Truth in the next noncumulative
language therefore defines \(\mathsf N_i(a)\) by
\(\exists y(\operatorname{Graph}_{i}(a,y)\land
\mathsf N_{i+1}(y))\), on every domain input. The named
atom is always a sentence, even when \(a\) is a noncode
and its old predicate value is false. Repeating this operation
lifts any lower noncumulative formula to a higher finite level.

Start the two desired translations as the identity at level zero.
Given the lower-level elementary graph \(u_{r-1}\), replace
a noncumulative \(\mathsf N_r(t)\) atom by the unique
graph output of \(u_{r-1}(t)\) followed by the cumulative
\(\mathsf C_r\) atom. On a valid input sentence, lower-level
value preservation gives the same truth; on an invalid input the
output is \(c(v_1)\), and both predicates are false. For
the reverse direction replace \(\mathsf C_r(t)\) by the
unique graph output of \(v_{r-1}(t)\) followed by
\(\mathsf N_r\). Lift every earlier cumulative predicate's
already supplied definition into the current noncumulative language
by the lifting just proved. These finitely many substitutions give
\(u_r,v_r\), whose code graphs are elementary by the annotated
certificate argument. No map invokes its own currently unconstructed
graph; only lower-level graphs occur in its fixed atom replacements.

Induction on the finite level and then on each actual finite formula
proves value preservation under every assignment. Ordinary atoms
are unchanged, truth atoms have the proved values, and the Boolean
constructors and quantifiers use the same domain and interpretation.
This proves the acceptable-structure finite comparison. It supplies
no unspecified transfinite limit convention or intrinsic-level
converse by fiat.
\end{proof}
\end{annexitem}

\begin{annexitem}{annex:p11.ktf.recursive-hierarchy}{thm:p11-ktf-recursive-hierarchy}{Existence of the completed classical hierarchy}
\begin{proof}
The ordinary finite alphabet and the numerical indices of the auxiliary
predicates have recursive tagged encodings. A finite word can be parsed,
its variables and binders checked, and its predicate indices listed
by a terminating numerical algorithm. The test for
\(\mathcal H_{<b}\) additionally requires every auxiliary
index \(d\) to be nonzero and to satisfy \(d\prec_W b\).
These are finitely many decidable tests using the fixed recursive
well-order. The analogous \(\mathcal H_b\) test uses
\(d\preceq_W b\). Thus the syntax predicates and usual
capture-free substitutions are computable uniformly in their numerical
parameters. Their arithmetic graphs follow from the independent
finite-history graph theorem, not from a truth classification.

Let \(\theta\) be the order type of \(W\), and use its
unique order isomorphism to identify each \(b\) with
\(\alpha_W(b)<\theta\). At zero assign no auxiliary
predicate. At any nonzero index suppose all strictly earlier predicate
extensions have been assigned. For each actual finite
\(\mathcal H_{<b}\)-formula and numerical assignment, recursion
on its finite expression gives its unique classical value: ordinary
atoms use \(\mathcal M\), and each auxiliary atom uses an
already assigned predicate with a strictly lower index. Quantifiers
range over the fixed set \(\mathbb N\). Separation on the
set of word codes gives the set of its true sentences. Use that
set for \(\mathsf{Tr}_b\), with false values on every other
number. This defines the next extension uniquely from the earlier
history; it makes no recursive request for \(\mathsf{Tr}_b\)
in its own input language.

For explicit existence, two partial histories satisfying these clauses
agree on their common initial segment, by a least-disagreement argument.
At a successor append the uniquely specified set just described.
At a limit take the union of the unique agreeing earlier histories;
Replacement collects this set of histories, and its union has the
required domain and clauses. Transfinite induction constructs the
whole history of length \(\theta\). The same least-disagreement
argument gives uniqueness. This is ordinary well-order recursion
on set-sized histories. Its output truth sets are not asserted
elementary merely because their syntax predicates are elementary.
It supplies the interpretation required by the definition for each
fixed \(\mathcal M\).
\end{proof}
\end{annexitem}

\begin{annexitem}{annex:p11.ktf.limit-separation}{thm:p11-ktf-limit-separation}{Completion is not a harmless renaming of a union}
\begin{proof}
Fix a nonzero \(b\) and its classical interpretation. Suppose
\(\chi(x)\in\mathcal H_{<b}\) defines exactly its true
sentence-code set, using the specified auxiliary coding. The diagonal
word map \(d_W(k)\) replaces the sole designated free variable
of the formula coded by \(k\) by \(\bar k\), and returns
a fixed noncode on invalid inputs. It is total computable, hence
has an arithmetic graph \(D_W(k,z)\) in the base language.
With fresh variables put
\[
 A(y)=\exists z\,(D_W(y,z)\land\neg\chi(z)),
 \qquad a=c_W(A),\qquad \ell=A(\bar a).
\]
Its only possible free variable before substitution is the designated
\(y\); choose that position as the diagonal map's variable,
or apply its computable renaming first. Then
\(d_W(a)=c_W(\ell)\). The unique graph output gives
\(\ell\leftrightarrow\neg\chi(c_W(\ell))\) in the
classical interpretation. All the predicates of \(\ell\)
are from \(\mathcal H_{<b}\): \(D_W\) is arithmetic
and \(\chi\) was assumed in that language. Therefore the
assumed truth definition also gives
\(\chi(c_W(\ell))\leftrightarrow\ell\), a contradiction
to classical bivalence. No infinitary formula or variable-indexed
predicate token was inserted into this diagonal sentence.

This proves undefinability of the full truth set in its own prior
language. At a limit that prior language is exactly the ordinary
union of the earlier completed languages, whereas
\(\mathcal H_b\) contains its genuine new truth predicate.
The two language conventions cannot be conflated on the ground
that they have the same preceding indexed predicates.
\end{proof}
\end{annexitem}

\begin{annexitem}{annex:p11.ktf.transfinite-embedding}{thm:p11-ktf-transfinite-embedding}{A uniform transfinite hierarchy embedding}
\begin{proof}
First construct a total word operation \(\operatorname{tr}(k,e)\)
without interpreting any truth predicate. On a valid auxiliary formula
word \(e\), copy each ordinary atom into the original coding
and preserve its ordinary terms. Replace an atom
\(\mathsf{Tr}_d(t)\) by a copy of the prospective original-language
template coded by \(k\), with its three designated free
variables simultaneously replaced by \(\bar k,\bar d,t\).
First move its bound variables and free parameter slots above
the largest variable indices in the input and template; then perform
the substitutions. This prevents both capture of variables in
\(t\) and accidental substitution into the newly inserted
term. The template is copied as a finite word, not recursively
translated again. Reject invalid templates or input words by returning
a fixed false base formula. Ordinary formula constructors are rebuilt
from their translated children. Parsing and these finite-word operations
terminate, so the operation is total computable. Its arithmetic
graph \(E(k,e,y)\) is supplied by the graph theorem.

Let \(\mathsf{Sent}_{<W}(b,e)\) be the arithmetic guard
for a \(\mathcal H_{<b}\)-sentence. Form the original
truth-language template
\[
 Q(k,b,x)=(b\in B)\land(b\ne0)\land
    \mathsf{Sent}_{<W}(b,x)\land
       \exists y\,(E(k,x,y)\land T(y)).
\]
It has only the three indicated free variables, chosen as the
operation's designated slots. Its ordinary guards and graph are
arithmetic and independent of the empirical predicates. Put
\(q=c_T(Q)\), define
\(\delta_W(F)\) to be the formula whose original word code
is \(\operatorname{tr}(q,c_W(F))\), and write
\(D_b(x)=Q(\bar q,\bar b,x)\).
In the translation each indexed atom is precisely \(D_d(t)\).
Inside \(D_b(e)\), the graph at a valid sentence code selects
exactly \(c_T(\delta_W(F))\). The constructor depended on
an arbitrary template code before \(Q\) was written; no uniform
truth-stage predicate or hierarchy interpretation was smuggled into
that word operation.

Fix \(\mathcal M\), and prove by well-order induction that
\(D_b\), for nonzero \(b\), has its prescribed classical
\(\mathsf{Tr}_b\)-values at every input by stage
\(\alpha_W(b)\). Invalid input sentence codes have a false
ordinary guard, so their translated predicate value is already
false at stage zero. On a valid sentence \(F\) of
\(\mathcal H_{<b}\), its finitely many indexed predicates
are \(\mathsf{Tr}_d\) with \(0\ne d\prec_W b\).
Their translations are \(D_d\), and the induction hypothesis
is uniform over all numerical arguments. Hence, at a stage containing
all their specified lower bounds, finite formula induction makes
\(\delta_W(F)\) bivalent with exactly its classical hierarchy
value under every assignment. Boolean constructors preserve these
values, and every quantifier instance is bivalent at that same
stage, since each indexed predicate's bound is uniform in its
argument. There is no illicit choice of a separate unbounded stage
for each quantifier instance.

If \(\alpha_W(b)=\gamma+1\), its predecessor index has
rank \(\gamma\). All finitely used lower predicate bounds
are at most \(\gamma\). Thus \(\delta_W(F)\) has its
classical value at \(S_{\mathcal M}^{\gamma}\), and its
code receives the corresponding truth-predicate sign at jump
\(\gamma+1=\alpha_W(b)\). The unique graph in \(D_b\)
then copies that sign. This includes rank one, where the sentence
is truth free at stage zero but its truth code is defined only
at jump one.

If \(\alpha_W(b)=\lambda\) is a nonzero limit, choose
\(\gamma\) to be the maximum of zero and the finitely
many ranks of the predicate indices in \(F\). Then
\(\gamma<\lambda\) and \(\gamma+1<\lambda\).
The translated sentence has its classical value at stage
\(\gamma\), and its code has its sign by
\(\gamma+1\), hence in the actual union stage \(\lambda\).
Every valid sentence input supplies such a finite list. Consequently
\(D_b\) is bivalent at all inputs by the single stage
\(\lambda\), even though no one smaller bound was asserted
for the whole input family. This proves the inductive claim at
limits without assuming that a truth jump commutes with their unions.

Now let \(F\in\mathcal H_b\). Every one of its indexed
predicates has index at most \(b\), so the preceding uniform
bounds and information persistence make every translated atom bivalent
at \(S_{\mathcal M}^{\alpha_W(b)}\). Formula induction
gives the stated equality for \(F\) under every assignment.
When \(b=0\) there are no indexed predicates and ordinary
formula translation gives its classical value at stage zero.
Defined stage values persist to the least fixed point and thence
to every fixed point extending it. The ordinary structure \(\mathcal M\)
was arbitrary, and all syntax graphs were fixed independently of
its empirical predicates, proving the uniform assertion over
\(\mathfrak B\).
\end{proof}
\end{annexitem}

\begin{annexitem}{annex:p11.ktf.stage-compiler}{thm:p11-ktf-stage-compiler}{Self-indexed finite syntax compilation}
\begin{proof}
An index is needed because the desired finite coordinate templates
refer to the graph of their own word compiler, not because a
truth-stage relation has already been defined. Given any numerical
prospective index \(i\), obtain an arithmetic formula
\(G_i(\sigma,d,x,y)\) for that program's computation graph.
The graph theorem makes this formula construction effective whether
or not the prospective program happens to be total. It is a
finite arithmetic word and contains no indexed truth tokens.

Define a total syntactic factory \(p(i,b)\) as follows. If
\(b=0\) or \(b\notin B\), output the formula \(x\ne x\).
Otherwise output the displayed word for \(\mathsf P_b(x)\)
with \(G_i\) as its graph formula and with the literal
predicate token \(\mathsf{Tr}_b\). Define \(n(i,b)\)
in the same way for the displayed negative formula. The well-order
guard has its fixed arithmetic word; fresh bound variable names
are chosen effectively. The two outputs belong to \(\mathcal H_b\)
for \(b\in B\) and have only their designated free
argument variable. The exceptional word \(x\ne x\) deliberately
retains that variable while being constantly false. This construction
runs no prospective program and evaluates none of its graph formula's
instances.

Using these factory words, construct a program \(f(i)\) which
takes \((\sigma,b,e)\) and transforms the finite original-language
formula coded by \(e\). At an ordinary atom, its positive
word is that atom and its negative word its negation. At a truth
atom \(T(t)\), substitute \(t\) capture freely into
\(p(i,b)\) and \(n(i,b)\), respectively. The remaining
paired recursions are explicitly
\[
\begin{array}{c|cc}
 F&F^+&F^-\\\hline
 \neg A&A^-&A^+\\
 A\land B&A^+\land B^+&A^-\lor B^-\\
 A\lor B&A^+\lor B^+&A^-\land B^-\\
 \exists x A&\exists x A^+&\forall x A^-\\
 \forall x A&\forall x A^+&\exists x A^- .
\end{array}
\]
Implication and other defined connectives are expanded first by
their fixed finite definitions. Capture-free copying, constructor
recursion and fresh-variable selection are finite word algorithms.
On invalid inputs output a fixed false base sentence. Thus
\(f(i)\) is a total computable program for every \(i\).
Constructing its index from \(i\) is total computable.

By the proved program self-reference lemma
\(\ref{thm:p04-doh-program-tools}\), take an index \(e\)
with \(\varphi_e=\varphi_{f(e)}\), keeping the numerical
arguments as program inputs. Because the right-hand program was
already total by finite syntax recursion, the actual program \(e\)
is total. Its formula \(G_e\) is therefore a unique-output
graph on every numerical input. Set \(\mathsf P_b=p(e,b)\),
\(\mathsf N_b=n(e,b)\), and let \(C_W^+,C_W^-\)
be its two actual outputs. This supplies the displayed coordinate
formulas coherently with their actual compiler index. On a valid
input they have only the auxiliary token \(\mathsf{Tr}_b\),
unless \(b=0\), and preserve free variables. No induction
on ordinal truth stages has been used in their existence or totality.

Finally observe the uniform value property of these syntactic
transformations. In any classical interpretation of the coordinate
formulas as two disjoint relations \(P,N\), the transformed
positive word is true exactly when the original formula has Strong-Kleene
value \(1\) with \(T\)'s coordinates \(P,N\), and
the negative word is true exactly when its value is \(0\).
This follows by finite formula induction using the displayed table
and the bivalence of ordinary atoms. In the truth-atom case
the factory supplies exactly the respective coordinate formula;
substitution is capture free. The universal negative condition
for an existential uses all instances and is exactly its Strong-Kleene
false condition. This value property is conditional on the supplied
coordinates; the next proof, rather than this syntax construction,
identifies them with the actual stages.
\end{proof}
\end{annexitem}

\begin{annexitem}{annex:p11.ktf.transfinite-stages}{thm:p11-ktf-transfinite-stages}{Exact coordinate representation through recursive transfinite levels}
\begin{proof}
Fix \(\mathcal M\in\mathfrak B\). Prove the coordinate
claim by well-order induction on \(b\). At \(b=0\)
both factory formulas are \(x\ne x\), so both extensions
are empty. This is the original stage zero, including its absence
of false nonsentence signs. Assume the claim at every
\(d\prec_W b\), with \(b\ne0\).

If \(x\) is an original sentence code, its two actual compiler
outputs at \(d\) are closed \(\mathcal H_d\)-formulas.
Their only possible indexed token is \(\mathsf{Tr}_d\),
or none when \(d=0\). Therefore they are sentences of
\(\mathcal H_{<b}\), where \(\mathsf{Tr}_b\) evaluates
their actual classical values. The unique-output graph \(G_e\)
selects precisely these words. The induction hypothesis identifies
the coordinate formulas at \(d\) with
\(S_{\mathcal M}^{\alpha_W(d)}\). The syntax compiler's
finite formula induction then identifies these two queried truth
values with the positive and negative sentence outputs of
\(J_{\mathrm K}(S_{\mathcal M}^{\alpha_W(d)})\).

The two displayed coordinate definitions existentially range over
all such lower \(d\). Their positive extension is consequently
the union of these lower positive jump outputs. Their negative
extension is the corresponding union of the negative jump outputs:
on original sentence codes it uses the negative compiler output;
on nonsentence codes its ordinary disjunct is true whenever there
is a lower index. Every nonzero \(b\) has one, since zero
is least. Thus these base false signs agree exactly with the
original first-jump convention. We have proved
\[
 (\mathsf P_b^{\mathcal M},\mathsf N_b^{\mathcal M})
    =\bigcup_{d\prec_W b}
          S_{\mathcal M}^{\alpha_W(d)+1},
\]
with coordinate unions. These are consistent increasing stages,
so the two computed relations are disjoint.

Write \(\alpha=\alpha_W(b)\). If \(\alpha=\gamma+1\),
the predecessor set has a largest index of rank \(\gamma\).
Its jump output is \(S_{\mathcal M}^{\gamma+1}\), and
all earlier outputs are included in it. If \(\alpha\) is
a nonzero limit, the ranks of lower indices are exactly the
ordinals \(\gamma<\alpha\). Hence
\[
 \bigcup_{\gamma<\alpha}S_{\mathcal M}^{\gamma+1}
   =\bigcup_{\gamma<\alpha}S_{\mathcal M}^{\gamma}
   =S_{\mathcal M}^{\alpha}.
\]
The first equality uses increasing stages and the fact that
\(\gamma+1<\alpha\) at a limit; the second is the actual
limit definition. It is not a jump-preservation assertion. This
completes the coordinate induction.

For arbitrary formulas and assignments apply the compiler's finite
value induction with these now-proved coordinate relations. The
positive and negative outputs test precisely the respective partial
values at \(S_{\mathcal M}^{\alpha_W(b)}\). A gap makes
neither condition true. The arithmetic graphs and word operations
are interpreted identically in every member of \(\mathfrak B\),
so the argument holds over the entire stated empirical family.
The auxiliary predicates were defined earlier as ordinary classical
truth on lower-indexed languages, not as the coordinates subsequently
proved equal to them through these formulas.
\end{proof}
\end{annexitem}

\begin{annexitem}{annex:p11.ktf.uniform-converse}{thm:p11-ktf-uniform-converse}{A proved uniform-level converse}
\begin{proof}
For a fixed interpretation and assignment, groundedness implies a
first defined evaluation stage: substitute the finitely many assigned
numerals into the free variables and use the evaluation and substitution
lemma. The bottom iteration for this countable sentence syntax stabilizes
at a countable stage, so the first stage is a countable ordinal.
For a safe formula these ranks form a set, since the empirical family
is a set and the relevant assignments are finite numerical tuples.
Their supremum \(\ell_{\mathfrak B}(F)\) is therefore an ordinal
at most \(\omega_1\). Each rank is at most that supremum.
Information persistence makes the value defined at the supremum for
every interpretation and assignment. Conversely, if one ordinal
\(\beta\) is a uniform bound, every first rank is at most
\(\beta\), so their supremum is at most \(\beta\).
This proves that the supremum is the least uniform ordinal bound.
It also proves the stated criterion for a represented bound, without
requiring a rank to attain the supremum or to be identical across
interpretations. Stage sequences are extended constantly after their
individual stabilizations, so the assertion at \(\omega_1\)
is well defined.

The exact-stage theorem constructs \(C_W^+(b,F)\) effectively
for every finite original-language formula, without deciding its
uniform-bound condition. In each \(\mathcal M\) its classical
value is true exactly when \(F\) is true at
\(S_{\mathcal M}^{\alpha_W(b)}\). Under the condition
that \(F\) is defined there for every interpretation and
assignment, the remaining case is actual falsehood, not a gap.
Thus this bivalent auxiliary formula equals \(F\)'s partial
value on every indicated interpretation and assignment. The negative
compiler output tests the complementary case under this same
condition. Information persistence preserves that defined value
to the least fixed point and to every fixed point above it,
giving the additional fixed-point equality.

Conversely apply the transfinite embedding to any
\(F\in\mathcal H_b\). Its one effective translation
\(\delta_W(F)\) has exactly the classical hierarchy value
at all interpretations and assignments, and is defined by the
uniform stage bound \(\alpha_W(b)\). This is the required
representative in the bounded partial-truth class. Composing the
compiler with the embedding produces a bounded original-language
representative of the same value, not merely an auxiliary formula
selected separately after discovering an actual interpretation's
truth value.

The two constructions establish expressive equivalence at each
represented level, including limit indices. They need not invert
one another as words, preserve a minimal syntactic vocabulary,
or decide uniform groundedness. The hypotheses explicitly fix
the empirical family, the recursive well-order and the completed
limit convention. This theorem consequently does not prove a
claim about every grounded sentence in one interpretation, or
about an unspecified canonical notion of intrinsic level over
all arbitrary acceptable structures and ordinals.

Finally, for any recursive ordinal \(\alpha\), a recursive
presentation of order type \(\alpha+1\) has an element of rank
\(\alpha\). Rename its least element to zero by a computable
permutation of the numerical domain. It supplies the required
\(W,b\). Conversely the order type of a recursive well-order,
and each of its initial segments at a fixed numerical element,
are recursive ordinals. By the definition of the least nonrecursive
ordinal they are below \(\omega_1^{\mathrm{CK}}\).
No assertion of a single recursive presentation cofinal in all
recursive ordinals follows from these separate presentations.
\end{proof}
\end{annexitem}

\begin{annexitem}{annex:p11.ktf.intrinsic-scope}{thm:p11-ktf-intrinsic-scope}{Safety is not the same as a recursive intrinsic bound}
\begin{proof}
Use the same effective finite-sequence coding as in the tree
construction, and one binary ordinary predicate \(R\). For each
prefix-closed well-founded tree \(\mathcal T\subseteq
\mathbb N^{<\omega}\) containing the empty sequence, interpret
\(R(s,t)\) as immediate childhood in that tree, false on all
other pairs. Let \(\mathfrak B_{\mathrm{wf}}\) contain exactly
these interpretations with standard arithmetic held fixed. This is
a nonempty set of structures, although the set of their empirical
predicate extensions is not assumed effective.

The total word operation which closes a prospective two-variable
template at its own code and a node numeral has an arithmetic
graph \(G(k,t,u)\), in the current coding \(c_T\). On invalid
template codes it returns zero. Construct the actual template
\[
 A(x_0,x_1)=(x_0=x_0)\land(x_1=x_1)\land
  \forall t\forall u\,
       ((R(x_1,t)\land G(x_0,t,u))\to T(u)),
 \qquad a=c_T(A),\qquad \eta_s=A(\bar a,\bar s).
\]
All bound variables in the fixed graph formula are first renamed
freshly. The closure operation acts on finite words independently
of \(R\)'s extension, and
\(G(a,t,u)\) holds exactly when \(u=c_T(\eta_t)\).
The two tautological conjuncts retain both template slots; in
particular different node numerals give different sentence strings.
Take \(\eta\) to be \(\eta_s\) at the fixed empty-root
sequence code. It is one sentence, not a model-dependent choice.

At any partial state, each genuine child contributes exactly the
truth value assigned to its sentence code; every graph-false or
child-false numerical instance of the implication is true. Hence
\(\eta_s\) is true exactly when all child codes are positive,
false exactly when one is negative, and undefined otherwise.
Induction on the bottom stages excludes a first negative family
code: a successor would require an earlier negative child code,
and a union limit cannot introduce a first sign. No family sentence
is therefore false. On a well-founded tree, induction on children
shows that its first true evaluation stage is
\[
 r_{\mathcal M}(\eta_s)=
   \sup\{r_{\mathcal M}(\eta_t)+1:R(s,t)\}.
\]
Indeed each child code enters one successor after its first true
evaluation. At the displayed supremum all those codes are present,
whereas at every lower stage at least one required code is absent.
The empty supremum gives true evaluation at stage zero. This is
exactly the tree-height recursion. Computability of the child
relation played no part in this ordinal induction: that relation
is here an ordinary bivalent empirical predicate, not an assumed
arithmetical definition of a nonrecursive order.

Every countable ordinal \(\alpha\) occurs as the root height
of a numerical tree. To see this without an effective-presentation
assumption, identify the elements of the countable well-order
\(\alpha\) with a subset of \(\mathbb N\), and take the
tree of finite strictly decreasing sequences in that order, together
with its empty root. At a nonempty node ending at ordinal
\(\beta\), well-founded induction gives height \(\beta\):
its possible next elements are exactly the smaller ordinals, with
supremum of their heights plus one equal to \(\beta\).
The empty root consequently has height
\(\sup_{\beta<\alpha}(\beta+1)=\alpha\); for
\(\alpha=0\) it is a leaf. The tree has no branch, since an
infinite decreasing ordinal sequence is impossible. Its numerical
coding need not be recursive, which is permitted for \(R\).
Thus \(\eta\) is grounded and true in every member of the
specified family, and its ranks range cofinally over the countable
ordinals. Their supremum is \(\omega_1\). In particular no
recursive ordinal is a uniform bound.

If the allowed family instead includes a tree with a branch through
the empty root, a first true stage for its root would yield a strictly
decreasing sequence of first ranks along that branch: each child code
must already be positive, so its evaluation rank is strictly smaller.
This is impossible. Falsehood was already excluded, so that root
sentence stays undefined. Restriction to well-founded empirical
interpretations is therefore a consequential assumption in the safe
example, not an assertion that this sentence is safe under every
possible binary predicate interpretation. Within the restricted
family it is also equivalent to \(0=0\); absence of a recursive
evaluation bound is not a proof of semantic nonrepresentability
in a finite Tarski language.

For the remaining finite distinctions, let \(R\) instead be one
unary predicate, allow all its interpretations, and put
\[
 F=R(0)\lor T(\overline{c_T(0=0)}).
\]
At stage zero its second disjunct is undefined. Thus its first
defined stage is zero when \(R(0)\) is true. When \(R(0)\)
is false, the sentence is initially undefined and becomes true
at stage one, when the code of \(0=0\) enters positively.
It is safe, has intrinsic bound one over this full family, and
is true at every least point. It is consequently equivalent to
\(0=0\), whose intrinsic bound is zero. This proves both
variation of individual ranks and failure of rank invariance
under the specified semantic equivalence.
\end{proof}
\end{annexitem}

\begin{annexitem}{annex:p11.ktf.semantic-order-syntax}{thm:p11-ktf-semantic-order-syntax}{Elementary opaque-index syntax and prospective-template compilers}
\begin{proof}
Start with the finite postorder certificate proved in
Lemma~\ref{thm:p11-ktf-internal-syntax}. For the raw auxiliary
language remove its \(T\)-prime case and add exactly one case:
an atomic code is \(\iota(8_{\mathcal A},d,t)\) with a
valid term code \(t\) and arbitrary payload \(d\in A\).
There are still finitely many constructor cases. Neither the payload
\(d\) nor a constant's payload is interpreted as a child of
the parse. Induction on the finite record indices gives actual finite
expressions and unique children exactly as in the original proof;
possible cyclic equations in the ambient tuple coding do not supply
a finite postorder certificate. The same annotated free-index lists
prove the formula and sentence guards.

For a lower-language guard at \(b\), require \(I(b)\),
and at every indexed prime record require \(I(d)\), \(d\ne o\)
and \(d\prec_I b\). For the nonstrict guard use
\(d\preceq_I b\) instead. Verification quantifies over the
internal finite record positions and checks the stated ordinary
relations. It is therefore elementary whether or not the index order
has a decision algorithm. Actual membership in the corresponding
language supplies the certificate and conversely the certificate's
finite induction establishes that membership. This proves the exact
legality predicates, not merely raw syntactic formation.

We need transformations with a variable prospective template code,
not just the fixed-template case previously proved. Here are the
additional finite certificates. Parse the input formula and the
prospective template. Compute their greatest variable indices by
the elementary numerical finite histories used in the syntax lemma.
Give each binder occurrence in a template copy a fresh positive
index above that maximum, according to its finite postorder position.
Replace its bound variable occurrences using the nearest matching
binder on their finite ancestor paths. Such a path is a finite
sequence of record indices, each adjacent pair verified by a child
pointer, so the proposed ancestor and nearest-binder checks are
first order over the finite certificate. Free occurrences are not
renamed. This gives a prescribed, hence unique, fresh version of
each template copy; its variable slots are kept as free slots.

Simultaneously replace the designated free slots by their indicated
terms. At a free variable record select the corresponding input
term code; at other term records copy the original code. Rebuild
the remaining constructors. The inserted terms are copied, not
substituted into again. The newly selected binder indices exceed
all their variable indices, so no inserted free variable is captured.
The name of a prospective code or an index is simply a constant
term code with that opaque domain payload. All checks use the
fixed tuple graphs, numerical finite lists and previously supplied
parse relations. Existence is actual finite traversal, and finite
induction proves uniqueness and the stated capture-free substitution.
This establishes a total elementary graph for copying a valid
prospective template with the required slot substitutions.

For \(E_I\), annotate the input raw auxiliary parse with one
output original code at every record. Require the prospective template's
free-index list to be exactly its three designated slots. Ordinary
primes are unchanged.
An indexed prime at payload \(d\) receives a copy of the
prospective three-slot template, with its slots replaced by
\(\dot k,\dot d,t\). Negation, conjunction and existential
records rebuild from their children's outputs. Each prime copy has
the finite certificate just described. No record asks for translation
of the inserted template's own subexpressions. Thus finite parse
induction gives exactly one output, with an elementary verifying
formula. For an invalid template or raw input send the output to
the fixed original false-sentence code. All tests of validity here
are raw syntactic tests, not calls to an index-order decision procedure.

For \(\mathsf C_I\), first verify that \(k\) is an
ordinary relation-template code with free indices exactly its five
designated slots. At \(b=o\) the two prime factory words are
the false condition \(x\ne x\). Otherwise form the displayed
\(p_I,n_I\) by fixed constructors: insert the literal indexed
token with payload \(b\), names of \(b,k\) and the selector
marker, and the actual copied prospective relation formula with
its remaining slots filled by the fresh bound variables and \(x\).
The words for \(I\), \(\prec_I\) and the original sentence
guard are fixed ordinary formulas. Their extensions are not evaluated
in this formation. Comparison with the one fixed object \(o\)
is an elementary equality case; no lower-index test is executed.

Annotate the original formula's postorder parse with two output
codes. Ordinary atoms contribute themselves and their negations;
a \(T(t)\)-prime contributes the two factory words with term
\(t\) substituted. Negation exchanges the children's outputs;
conjunction uses conjunction of the positives and disjunction of
the negatives; existential uses existential of the positive and
universal of the negative. These are precisely the earlier signed
constructor table, with defined connectives expanded by their finite
definitions. Copying a factory does not recompile its contained
prospective relation formula. The annotation conditions and all
substitution certificates are first order over actual finite lists,
so the relation between \(k,\sigma,b,e\) and its selected
output \(y\) is elementary. Invalid selector, template or formula
inputs have the fixed auxiliary false-word output. The finite
traversal supplies an output on every tuple of domain objects,
and finite induction proves uniqueness.

Both defining graph formulas are ordinary formulas with finitely
many fixed coding parameters. Their existence used only the supplied
acceptable scheme and finite certificates; the induction-on-internal-
numbers resource is the one proved from that scheme, agreeing with
the original 5A.1 resource \cite[pp.~65--66]{Moschovakis1974EIASOriginal}.
No prospective formula was assumed to define this graph, no ordinal
stage was queried, and no uniform classical valuation predicate was
introduced. Finally, for any supplied disjoint interpretations of
the two coordinate factory formulas, finite formula induction with
the signed table identifies the two output conditions with exactly
the Strong-Kleene true and false values. The next coordinate theorem
establishes which stage coordinates those factory formulas denote.
\end{proof}
\end{annexitem}

\begin{annexitem}{annex:p11.ktf.semantic-order-embedding}{thm:p11-ktf-semantic-order-embedding}{Completed hierarchies over semantic well-orders}
\begin{proof}
The index domain is a set, so its well-order has a set order type
\(\theta\). Every actual auxiliary expression is finite; its
opaque constants and indexed primes are coded by the fixed tuple
scheme. Thus each language's sentence codes form a subset of \(A\).
At an index \(b\), given the strictly earlier predicate
extensions, classical recursion on these actual finite formulas
gives each sentence its unique value. Quantifiers range over the
fixed set \(A\). Separation produces its true-code subset of
\(A\); assign exactly that set to \(\mathsf{Tr}^{I}_b\).
This step makes no request for the predicate at its own index in
the prior language. At a limit the agreeing earlier histories have
a set union, collected by Replacement. Ordinary well-order recursion
therefore supplies a whole history, and induction on the least
disagreement proves uniqueness. All guards retain their previously
proved ordinary meanings. Nothing requires recursive enumeration of
the order or numerical representation of its index payloads.

For the undefinability assertion, suppose an auxiliary formula
\(\chi(x)\in\mathcal J_{<b}\) defined the full true-sentence
set of that same language. The raw auxiliary syntax proof gives an
elementary graph \(D_I(k,z)\) for closing a formula whose only
possible free variable is \(v_1\) at the name of its own code,
with a default on invalid
inputs. In the ordinary language this uses only finite syntax and
the held coding resources. With fresh variables form
\[
 B(v_1)=\exists z(D_I(v_1,z)\land\neg\chi(z)),\qquad
 a=c_I(B),\qquad \lambda_I=B(\dot a).
\]
The unique closing output is \(c_I(\lambda_I)\), so the
classical value of \(\lambda_I\) is the negation of
\(\chi(c_I(\lambda_I))\). Its finitely many indexed tokens
are exactly those of \(\chi\), hence are all lower than \(b\).
The assumed truth definition gives the opposite equality with that
same sentence value, contradicting classical bivalence. This is an
actual finite expression; it uses no variable-indexed truth primitive.

Now use the already constructed elementary graph \(E_I\) in
the original-language template
\[
 Q_I(v_1,v_2,v_3)=I(v_2)\land(v_2\ne\dot o)\land
  \mathsf{Sent}^{I}_{<}(v_2,v_3)\land
  \exists y(E_I(v_1,v_3,y)\land T(y)).
\]
Let \(q_I=c(Q_I)\), take \(\Delta_I(F)\) to be the
unique output formula of \(E_I(q_I,c_I(F),y)\), and write
\(D_b(x)=Q_I(\dot q_I,\dot b,x)\). Each translated indexed
prime is precisely a capture-free instance of \(D_d\) on its
original term. The word transformation preceded this choice of its
template code, so this is literal finite closing, not an assumed
semantic satisfaction relation.

Fix a member \(\mathcal M\) of the held family. Well-order
induction proves that every nonzero \(D_b\) has its intended
classical \(\mathsf{Tr}^{I}_b\) value at every domain input
by stage \(\beta_I(b)\). An input outside the lower-language
sentence set has a false ordinary guard, giving false value already
at stage zero. A valid input codes a finite sentence \(F\)
whose indexed tokens have indices \(d\prec_I b\). Their
translations have, by induction, the prescribed values uniformly
over every domain argument by their respective stages \(\beta_I(d)\).
Finite formula induction consequently makes \(\Delta_I(F)\)
bivalent, with its classical value, at any stage above the finitely
many such bounds. This uniformity over each indexed predicate's
arguments is essential: all instances of each quantifier over \(A\)
are defined at that same bound.

If \(\beta_I(b)=\gamma+1\), the strict predecessor segment
has a largest index of rank \(\gamma\). All tokens used in
\(F\) have rank at most \(\gamma\); hence its translated
sentence has its classical value at \(S_{\mathcal M}^{\gamma}\).
Its original sentence code enters with the corresponding sign at
the next jump, giving the correct \(D_b\) value at
\(\gamma+1\). If \(\beta_I(b)=\lambda\) is a nonzero
limit, let \(\gamma\) be the maximum of zero and the finitely
many used token ranks. Then \(\gamma<\lambda\) and
\(\gamma+1<\lambda\). The translated sentence is defined
at \(\gamma\), and its code has entered by \(\gamma+1\),
hence is present in the actual union stage \(\lambda\).
Every valid sentence input admits this finite maximum. The whole
\(D_b\) predicate is therefore defined at every input by the
one limit stage, without asserting a single smaller bound for all
its inputs or preservation of unions by the truth jump.

For any \(F\in\mathcal J_b\), its finitely many indexed
primes now have their uniform intended values by \(\beta_I(b)\).
Ordinary atoms are bivalent from the start. Formula induction yields
the claimed translation equality under every assignment. At \(b=o\)
the language has no indexed token, so the translation is the ordinary
constructor identity and has its value at stage zero. Information
persistence carries these values to the least fixed point and to
every other fixed point above it. The syntax graphs and the fixed
order are shared throughout \(\mathfrak F\), so the same proof
applies to every allowed ordinary interpretation.
\end{proof}
\end{annexitem}

\begin{annexitem}{annex:p11.ktf.semantic-order-stages}{thm:p11-ktf-semantic-order-stages}{Exact coordinates from a named compiler graph}
\begin{proof}
Choose a concrete ordinary defining formula for the total graph
\(\mathsf C_I(k,\sigma,b,e,y)\), renaming its designated
free slots to the factory's five positions. Conjoin their self-equalities
if a slot would otherwise be absent. It has only finitely many
fixed coding parameters, all of which have names, and no truth token.
Let its actual code be \(h\). The factory at \(k=h\) copies
this formula, inserts \(\dot h\) into its first slot and the
appropriate selector name into its second. Therefore the copied
relation is exactly the actual compiler graph at those arguments.
The definition of that graph required only verification of finite
copying and compilation certificates for an arbitrary prospective
code. It neither contained the then-unchosen \(h\) as a fixed
parameter nor depended on any stage interpretation. Naming it now
creates no infinite word: code objects inside constant terms are
opaque payloads, and the inserted graph copy is not recursively
compiled. This establishes the actual finite \(P^I_b,N^I_b\)
formulas and total signed outputs specified in the statement.

For any valid \(b\), the factory's only indexed token is
\(\mathsf{Tr}^{I}_b\), or none when \(b=o\). Its other
formulas, including the copied compiler graph, are ordinary. Hence
the actual signed outputs belong to \(\mathcal J_b\) and
preserve the input's free variables. The finite conditional value
induction proved in the syntax lemma applies once the factory
coordinate extensions are known to be disjoint.

Identify those extensions by well-order induction. At \(b=o\)
both are \(x\ne x\), giving exactly the original empty stage.
At a nonzero index, an original sentence code \(x\) has at
every \(d\prec_I b\) two actual closed compiler outputs.
They are sentences of \(\mathcal J_d\), hence of the prior
language \(\mathcal J_{<b}\). The graph occurrence selects
exactly those output codes, and \(\mathsf{Tr}^{I}_b\) has
their actual classical values by the hierarchy theorem. The
inductive coordinate identity at \(d\), followed by the finite
signed value induction, identifies these truth queries with the
positive and negative sentence outputs of
\(J_{\mathrm K}(S_{\mathcal M}^{\beta_I(d)})\).
For a nonsentence object the positive guard excludes membership;
the negative factory includes it whenever a lower index exists.
Every nonzero index has \(o\) below it. Thus negative noncodes
are supplied at exactly the same first-jump convention as in the
original stage sequence.

The two resulting extensions are consequently
\[
 (P^I_b{}^{\mathcal M},N^I_b{}^{\mathcal M})
       =\bigcup_{d\prec_I b}S_{\mathcal M}^{\beta_I(d)+1}.
\]
At a successor rank the largest preceding output is exactly that
stage. At a nonzero limit rank \(\lambda\), the lower index
ranks are precisely all \(\gamma<\lambda\), and the
increasing-stage identity gives
\[
 \bigcup_{\gamma<\lambda}S_{\mathcal M}^{\gamma+1}
   =\bigcup_{\gamma<\lambda}S_{\mathcal M}^{\gamma}
   =S_{\mathcal M}^{\lambda}.
\]
These are actual coordinate unions. They are disjoint, so the
conditional signed value induction is justified at this index as
well. Induction gives the exact coordinates at every index; then
finite formula induction under every assignment gives exactly the
two sign tests, with neither condition true for a gap. At no point
was a relation meaning stage membership assumed in order to obtain
the compiler graph or its named template.
\end{proof}
\end{annexitem}

\begin{annexitem}{annex:p11.ktf.semantic-order-converse}{thm:p11-ktf-semantic-order-converse}{Uniform converse and the actual ordinal reach}
\begin{proof}
Well-order the signed domain \(A\times\{0,1\}\), whose
cardinality is \(\kappa\). If the original increasing stage
sequence failed to stabilize before \(\kappa^+\), every
\(\alpha<\kappa^+\) would have a new sign in
\(S^{\alpha+1}\setminus S^\alpha\). Choose its least
signed domain object in the held well-order. Increasing stages make
these selected objects distinct for distinct \(\alpha\),
giving an injection of \(\kappa^+\) into a set of cardinality
\(\kappa\), a contradiction. Thus each fixed interpretation
stabilizes at some \(\xi<\kappa^+\). This argument counts
actual truth-coordinate signs, including first-jump noncodes; it
does not invoke the earlier countable-sentence bound in this regime.

A safe formula under a finite domain assignment has a first defined
stage below \(\kappa^+\): close it at the assignment's named
objects using the elementary closing operation, and use substitution
and the just-proved stabilization bound. Across the set family
\(\mathfrak F\) and its finite domain assignments these ranks
form a set. Their supremum \(\ell_{\mathfrak F}(F)\) is
at most \(\kappa^+\). Each first rank is at most the
supremum, so information persistence defines each value at the
supremum even if no instance attains it. Conversely any common
ordinal bound is at least every rank, hence at least their supremum.
This proves the stated index-bound criterion for the fixed order.
Sequences are extended constantly after their individual stabilizations.

At a represented bound \(b\), the exact-stage compiler's
positive condition is true exactly in the true case. Safety and
the common bound exclude the gap case, so its classical false
case agrees exactly with actual falsehood. It is therefore an
equivalent \(\mathcal J_b\)-formula over every interpretation
and assignment. The hierarchy embedding supplies the reverse
direction with the same uniform index bound. Defined stage values
persist at every fixed point over each interpretation. These are
the same semantic equivalence criteria as previously specified,
with finite domain rather than numerical assignments; no assertion
about minimum expressive complexity follows.

Because \(I\subseteq A\), its order type \(\theta\)
has cardinality at most \(\kappa\) and is strictly below
the initial ordinal \(\kappa^+\). Each index rank has that
same upper limitation. Conversely for any \(\alpha<\kappa^+\),
the well-order \(\alpha+1\) has cardinality at most \(\kappa\)
and can be injected into \(A\), taking its first element to
the designated \(o\). Transport its order to its image. If
that order is not elementary in the old structure, add precisely
the indicated unary index-set and binary order predicates as an
ordinary finite expansion. The old elementary coding scheme
remains elementary there. This supplies an actual comparison in
that expanded language for an index of rank \(\alpha\);
it is not a definability theorem for that order in the old language.
For a numerical domain these are exactly the countable ordinals,
with no recursive-order limitation. A set-sized fixed code domain
cannot name more distinct index tokens than its cardinality, so
this construction supplies no rank \(\kappa^+\) or all-ordinal
hierarchy on that same domain.

For the asserted direct finite-word arithmetic version, keep the
previous \(0,s,+,\times\) term alphabet and finite empirical
predicate vocabulary, adjoining the ordinary \(I,\prec_I\)
tokens and numerical indexed auxiliary tokens. Parse raw words
with these finitely many ordinary term-constructor cases and one
numerical indexed-prime case. Free-index lists, simultaneous
capture-free substitutions, fresh renamings and constructor copying
are the same terminating finite-word algorithms as before; ordinary
function terms are copied or substituted with their full finite
term trees. Names of numerical codes and indices are the actual
finite numerals. The prospective-template compilers build those
numerals and copy the finite order-guard words without querying
the order extension. They are total computable transformations
on raw word codes even when that extension is nonrecursive, so
Theorem~\ref{thm:p04-doh-arithmetic-graphs} supplies their actual
ordinary arithmetic graph formulas. Choose the actual word code
of the signed graph formula for its prospective template exactly
as in the opaque proof. No program-index fixed point is needed;
inserting the finite numeral of that code does not recompile
the copied graph formula. Language legality uses its ordinary
\(I,\prec_I\) guards. The hierarchy, embedding and coordinate
inductions then apply verbatim to these actual words: the only
extra atomic cases are bivalent ordinary function-term atoms.
Quantifiers still range over standard \(\mathbb N\), and
the fixed empirical order is a genuine countable well-order.
Thus this is a direct extension of the earlier arithmetic syntax,
not an unstated identification with a different relational coding.

If ordinary interpretations also vary the index-set or order
extensions, the syntax transformations still execute no semantic
order tests. Each structure satisfying the well-order hypothesis
has its own unique hierarchy, and the two induction proofs apply
using its own index ranks. Nothing then identifies the ordinals
assigned to a common opaque index across structures. This is why
the single-ordinal supremum criterion in the statement holds the
order fixed. The added resources and this family distinction are
essential hypotheses, not consequences of naming the intended
historical intrinsic-level interpretation.
\end{proof}
\end{annexitem}

\begin{annexitem}{annex:p11.ktf.expansion-transport}{thm:p11-ktf-expansion-transport}{Stage-preserving recoding into an ordinary expansion}
\begin{proof}
There is a genuine reason to recode rather than keep source strings.
An old truth atom can name an object that is an old nonsentence but
the code of a new target sentence. The old interpretation makes
that atom false after its first jump, whereas the target may make
it true or undefined. Enlarging the ordinary vocabulary therefore
does not by itself prove stage preservation for unchanged formulas.
For example, when \(I\) is newly adjoined, the object
\(c_1(I(\dot o))\) is an old nonsentence code but a new
truth-free true sentence code. The truth atom at that object is
false at the old stage one and true at the new stage one.

Use the same prospective-template certificate mechanism, with a
source parse and a target parse. Given an ordinary target template
code \(k\) with exactly three designated indices free, transform
a valid source formula by preserving its ordinary atoms, terms and
constructors, and replacing every source \(T(t)\) by
\(\exists y(K_k(k,t,y)\land T(y))\). Here \(K_k\)
is the actual capture-free copy of that template, with \(\dot k\)
in the first slot and \(t,y\) in the remaining slots. Its
new binder is fresh. The inserted copy is not transformed again.
Postorder annotation and the simultaneous-substitution certificates
from Lemma~\ref{thm:p11-ktf-semantic-order-syntax} prove existence,
uniqueness and an elementary graph for this finite operation, uniformly
in \(k\). Old ordinary primes still name their old predicates;
the additional target predicate cases occur only inside prospective
template copies. Variable and constant codes use the held scheme.

For an invalid source formula code, give the output
\(a_* = c_1(v_1)\). This is a variable-term code, not a formula
or sentence code: its constructor tag is different from every target
formula constructor, including the additional ordinary predicates.
No assumption that a numerical marker is a noncode is needed.
An invalid prospective template also has this specified output.
Call the total elementary graph of the operation \(\mathsf R(k,e,y)\).
Choose its actual ordinary defining formula, retain all three free
slots by self-equalities if necessary, and set its target code to
\(q_R\). The copied template at \(k=q_R\) is then precisely
that actual graph with its first slot named. Let \(\Lambda(F)\)
be the unique output formula on \((q_R,c_0(F))\), and let
\(\mu_R(a)\) be the unique output object on \((q_R,a)\).
This names the graph after its construction; it neither assumes
the desired stage equality nor executes the semantic formula it copies.

For every valid source formula, finite induction proves free-variable
preservation and the claimed translation shape. In particular a source
sentence maps to the target sentence \(\Lambda(F)\). A valid
source open formula maps to an open target formula and hence a target
nonsentence. An invalid source formula maps to the term code \(a_*\).
Thus for every \(a\in A\),
\[
 a\text{ is a source sentence code}\quad\Longleftrightarrow\quad
 \mu_R(a)\text{ is a target sentence code},
\]
and in the sentence case
\(\mu_R(c_0(F))=c_1(\Lambda(F))\). The equivalence uses
retention of the truth-prime term variable in the graph template;
it is not a claim that an arbitrary constant-output transformation
preserves sentencehood.

Prove by ordinal induction that the source truth atom on every
\(a\) has the same sign as the target truth atom on \(\mu_R(a)\).
At stage zero both truth predicates are wholly undefined. At any
stage where this atom assertion holds, finite formula induction gives
the asserted equality for \(F\) and \(\Lambda(F)\): ordinary
atoms have the same old interpretation; the graph in a translated
truth prime is bivalent with exactly one output, so its existential
has exactly that selected truth-atom value. False graph instances
cannot create a true witness or prevent a uniformly false existential.
Negation, conjunction and both quantifier signs then agree under
every assignment on the same domain.

At the next successor, if \(a=c_0(F)\) is a source sentence
code, its new sign is the preceding value of \(F\). The selected
target code is \(c_1(\Lambda(F))\), whose new sign is the
preceding value of \(\Lambda(F)\), equal by the finite induction.
If \(a\) is a source nonsentence, the selected object is a
target nonsentence and both jump outputs assign it false, regardless
of the previously added predicates. At a limit, membership of each
fixed \(a\) or \(\mu_R(a)\) in either coordinate is membership
at some earlier stage. The inductive atom equalities identify these
memberships through the two actual coordinate unions. This completes
the stage induction and the simultaneous formula-value equality.

The two stage sequences may stabilize at different ordinals because
the target has additional sentences. Extend each constantly after its
stabilization and compare them at a stage above both. Their least-point
formula values agree, and equality at every earlier stage gives exactly
the same first defined rank and groundedness. For a held family of
ordinary expansions this proof applies uniformly with the same graph,
as long as the old predicates and coding data agree as stipulated.
The semantic-order converse can now be applied to the actual translated
formula, not to an unjustified unchanged string, giving the final
source-to-hierarchy conclusion at its common represented bound.

In the numerical word-coding case, the analogous source/target parses
and copied-template substitutions are total finite-word algorithms.
Use a fixed invalid term word as the nonsentence default, and obtain
the arithmetic graph by the independent graph theorem. Finite function
terms are preserved as in the direct arithmetic proof above. The same
stage induction therefore proves the numerical assertion as well.
No inverse translation of an arbitrary added ordinary predicate was
constructed; without its old-base definition such an assertion would
require an additional premise.
\end{proof}
\end{annexitem}
